
\subsection[Category of Template Instantiations]{The category of template instantiations $\tinstC$}

We extend category $\tempC$ with instantiation data. 
This is where our template heap functor will be indexed over.
\begin{itemize}
    \item Objects: Sort-preserving injective renamings between worlds $\Minus\vw \injto \Plus\vw$.
    \item Morphisms: Injective, sort-preserving $\concL\me$ and $\anyL\me$ such that
                    \begin{center}
                    \begin{tikzcd}{\Plus\vw} \ar[r, tail, "{\anyL\me}"] & {\Plus{\hat\vw}} \\
                                  {\Minus\vw} \ar[u, tail, "\sigma"] \ar[r, tail, "{\concL\me}"'] & {\Minus{\hat\vw}} \ar[u, tail, "\hat\sigma"']
                     \end{tikzcd}
                    \end{center}
                    and instantiation data $\mvarM_\me : \dheapF(\Minus{\hat\vw} \ominus \concL\me, \Plus{\hat\vw})$.
\end{itemize}
We have a forgetful functor $\tinstU : \tinstC \to \tempC$, and an inclusion functor
$\inclITI\dash : \instC \inclto \tinstC$ which maps to its identity morphism.
We inherit the terminology of $\tempC$ and talk about concrete locations and location variables, 
we overload functors $\concL\dash,\anyL\dash : \tinstC \to \worldC$ by precomposing them with $\tinstU$,
and talk about variable preserving and reflecting morphisms by noticing that $\tinstU$ reflects these properties.

\begin{definition}[Composition of morphisms in $\tinstC$]
Composition of $\sigma_1 \xto{\me_1} \sigma_2 \xto{\me_2} \sigma_3$ is defined as composition of the underlying injections and composition of the instantiation data is given by
\begin{align*}
        \mvarM_{\me_2 \comp \me_1} &: \dheapF(\concL{\sigma_3} \ominus \concL{\me_2 \comp \me_1}, \anyL{\sigma_3}) \\
        \va : \vC & \mapsto \begin{cases}\sem{\phi(\vC)}\anyL{\me_2}(\proj_{\va'}\mvarM_{\me_1}) & \mif \concL{\me_2}(\va') = \va \\
                                         \proj_\va{\mvarM_{\me_2}} & \mif \va \in \concL{\sigma_3} \ominus \concL{\me_2}
                            \end{cases}
\end{align*}
\end{definition}

\begin{remark}
    Every $\sigma, \mvarM \in \dheapF(\concL\sigma, \anyL\sigma)$ induces a morphism $?_\mvarM$
    \begin{center}
        \begin{tikzcd}
            {\emptyset} \ar[rr, "?_{\anyL\sigma}"] && {\anyL\sigma} \\
            {\emptyset} \ar[u, "?_\emptyset"] \ar[rr, "?_{\concL\sigma}"'] && {\concL\sigma} \ar[u, "\concL\sigma"']
        \end{tikzcd}
    \end{center}
\end{remark}

\begin{definition}[Template Heap Functor]
We define the template heap functor $\theapF : \tinstC \to \setC$ as
\begin{align*}
    \theapF(\sigma_1 \xto\me \sigma_2) &: \dheapF(\concL{\sigma_1}, \anyL{\sigma_1}) \to \dheapF(\concL{\sigma_2}, \anyL{\sigma_2}) \\
    \theapF(\me)(\mvarM') &= \mvarM_{\me \comp ?_{\mvarM'}}
\end{align*}
Note that $\mvarM_\me : \theapF(\concL{\sigma_2} \ominus \concL\me \xto{\compl{\concL{\me}}} \concL{\sigma_2} \xto{\concL{\sigma_2}} \anyL{\sigma_2})$ 
and we have that $\theapF \iso \homF\tinstC{?_\emptyset}\dash$.
\end{definition}

\begin{definition}[Variable Extension Functor]
We lift the variable extension functor from $\tempC$ to $\tinstC$
\[\begin{array}{r@{\,}c@{\,}l}
    \extF\dash\vw & : & \tinstC \to \tinstC \\
    \extF\sigma\vw & = & \incl_1 \comp \concL\sigma \\
    \extF\me\vw & = & \Pair{\extF{\tinstU\me}\vw}{\dheapF\pair1{\incl_1}(\mvarM_\me)}
\end{array}\]
\end{definition}

\begin{lemma}[Inhabitation in $\tinstC$]\label{def:inhabited-sorts-tinstC}
    A set of sorts is inhabited (\cref{def:inhabited-sorts}) iff
    for every template $\sigma$, there exists a concrete template $\inclITI{\vw'}$ and a morphism $\sigma \xto\me \inclITI{\vw'}$.
\end{lemma}
\begin{proof}
    For the $\to$ direction, consider $\varL\sigma$ and so for every $\va : \vC \in \varL\sigma$ we have a $\va : \vC \to \vw$.
    By induction on world $\varL\sigma$, this gives us a $\varL\sigma \to \hat\vw$ and a heap $\mvarM : \heapF(\hat\vw)$ by taking the disjoint sum of the worlds and heaps.
    We use this data to construct a morphism from $\sigma$ to $1_{\concL\sigma \oplus \hat\vw}$.

    For the $\from$ direction, we have a morphism $?_{\va : \vC} \xto\me 1_\vw$ for every $\va : \vC$ which is exactly the data needed.
\end{proof}


\begin{lemma}[Monomorphisms in $\tinstC$]\label{lem:monos-tinistC}
    Every morphism in $\tinstC$ is mono.
\end{lemma}
\begin{proof}
    We show $\hat\me : \sigma_2 \to \hat\sigma$ is mono.
    Consider $\hat\me \comp \me_1 = \hat\me \comp \me_2$. 
    We have that $\tinstU\me_1 = \tinstU\me_2$ from \cref{lem:monos-T}.
    It remains to show $\mvarM_{\me_1} = \mvarM_{\me_2}$, so take $\va : \vC \in \concL{\sigma_2} \ominus \concL{\me_1}$.
    \begin{align*}
        \proj_\va\mvarM_{\me_1} & = \proj_\va\mvarM_{\me_2} \iff \\
        \sem{\phi(\vC)}\anyL{\hat\me}(\proj_\va\mvarM_{\me_1}) & = \sem{\phi(\vC)}\anyL{\hat\me}(\proj_\va\mvarM_{\me_2}) \iff \\
        \proj_{\concL{\hat\me}(\va)}(\mvarM_{\me \comp \me_1}) & = \proj_{\concL{\hat\me}(\va)}(\mvarM_{\me \comp \me_2})
    \end{align*}
\end{proof}

\begin{lemma}[Forgetful Functor Preserves Equalizers]\label{lem:tinstU-eq-pres}
    $\tinstU : \tinstC \to \tempC$ preserves equalizers. 
\end{lemma}
\begin{proof}
    Let $\ve : \sigma_1 \to \sigma_2$ be an equalizer of $\me_1$, $\me_2$.
    We show that $\sigma_1$ is canonically isomorphic to $\eqz{\tinstU{\me_1}}{\tinstU{\me_2}}$ when extended with appropriate instantiation data.

    \begin{center}
        \begin{tikzcd}
            {\eqz{\tinstU{\me_1}}{\tinstU{\me_2}}} 
                \ar[rr, "{\eqz{\tinstU{\me_1}}{\tinstU{\me_2}}}"]
                && {\sigma_2} 
                    \ar[r, yshift=3pt, "\tinstU{\me_1}"] \ar[r, yshift=-3pt, "\tinstU{\me_2}"']
                & {\sigma_3}
            \\
            {\sigma_1} 
                \ar[urr, "\tinstU\ve"'] 
                \ar[u, dashed, "\vh"]
        \end{tikzcd}
    \end{center}
    We equip $\vh$ and $\eqz{\tinstU\me_1}{\tinstU\me_2}$ by splitting $\mvarM_\ve$ using the canonical isomorphism
    \[ \concL{\sigma_2} \ominus \concL{\tinstU\me} \iso (\concL{\sigma_2} \ominus \eqz{\concL{\tinstU\me_1}}{\concL{\tinstU\me_2}}) \oplus
                                                        (\eqz{\concL{\tinstU\me_1}}{\concL{\tinstU\me_2}} \ominus \concL\vh) \]
    and we show that the $(\eqz{\concL{\tinstU\me_1}}{\concL{\tinstU\me_2}} \ominus \concL\vh)$ part of data $\mvarM_\ve$ has locations in $\eqz{\anyL{\tinstU\me_1}}{\anyL{\tinstU\me_2}}$.

    So consider $\va : \vC \in \eqz{\concL{\tinstU\me_1}}{\concL{\tinstU\me_2}} \ominus \concL\vh$, and $\pi_{\eqz{\concL{\tinstU\me_1}}{\concL{\tinstU\me_2}}(\va)} \mvarM_\ve$. 
    Notice that
    \begin{align*}
        \mvarM_{\me_1 \comp \ve} &= \mvarM_{\me_2 \comp \ve} \implies  & \tag{def of $\eqz{\concL{\tinstU\me_1}}{\concL{\tinstU\me_2}}$} \\
        \pi_{\concL{\me_1}(\eqz{\concL{\tinstU\me_1}}{\concL{\tinstU\me_2}}(\va))} \mvarM_{\me_1 \comp \ve} & = 
        \pi_{\concL{\me_2}(\eqz{\concL{\tinstU\me_1}}{\concL{\tinstU\me_2}}(\va))} \mvarM_{\me_2 \comp \ve} \implies      & \tag{def of $\mvarM_{\me' \comp \me}$} \\
        \sem{\phi(\vC)}\anyL{\me_1}(\pi_{\eqz{\concL{\tinstU\me_1}}{\concL{\tinstU\me_2}}(\va)} \mvarM_\ve) & = 
        \sem{\phi(\vC)}\anyL{\me_2}(\pi_{\eqz{\concL{\tinstU\me_1}}{\concL{\tinstU\me_2}}(\va)} \mvarM_\ve)
    \end{align*}
    and so for every $\vb : \vC' \in \anyL{\sigma_2}$, we have 
    \[ \vb \isIn \pi_{\eqz{\concL{\tinstU\me_1}}{\concL{\tinstU\me_2}}(\va)} \mvarM_\ve \implies \anyL{\me_1}(\vb) = \anyL{\me_2}(\vb) \]

    We show that $\mvarM_{\me_1 \comp \eqz{\tinstU\me_1}{\tinstU\me_2}} = \mvarM_{\me_2 \comp \eqz{\tinstU\me_1}{\tinstU\me_2}}$
    and so consider $\va : \vC \in \concL{\sigma_3} \ominus \concL{\me_1 \comp \eqz{\tinstU\me_1}{\tinstU\me_2}}$ and proceed by cases
    \begin{itemize}
        \item $\va \in \concL{\sigma_3} \ominus \concL{\tinstU\me_1}$ and $\va \in \concL{\sigma_3} \ominus \concL{\tinstU\me_2}$)
            \begin{align*}
                \proj_\va \mvarM_{\me_1 \comp \eqz{\tinstU\me_1}{\tinstU\me_2}} & = \proj_\va \mvarM_{\me_2 \comp \eqz{\tinstU\me_1}{\tinstU\me_2}} \iff \\
                \proj_\va \mvarM_{\me_1} & = \proj_\va \mvarM_{\me_2} \iff \\
                \proj_\va \mvarM_{\me_1 \comp \ve} & = \proj_\va \mvarM_{\me_2 \comp \ve}
            \end{align*}
        \item $\va \in \concL{\sigma_3} \ominus \concL{\tinstU\me_1}$, $\va' \in \concL{\sigma_2} \ominus \eqz{\concL{\tinstU\me_1}}{\concL{\tinstU\me_2}}$,
              $\va = \concL{\me_2}(\va')$)
            \begin{align*}
                \proj_\va \mvarM_{\me_1 \comp \eqz{\tinstU\me_1}{\tinstU\me_2}} & = \proj_\va \mvarM_{\me_2 \comp \eqz{\tinstU\me_1}{\tinstU\me_2}} \iff \\
                \proj_\va \mvarM_{\me_1} & = \sem{\phi(\vC)}\anyL{\me_2}(\proj_{\va'} \mvarM_{\eqz{\tinstU\me_1}{\tinstU\me_2}}) \iff \\
                \proj_\va \mvarM_{\me_1} & = \sem{\phi(\vC)}\anyL{\me_2}(\proj_{\va'} \mvarM_\ve) \iff \\
                \proj_\va \mvarM_{\me_1 \comp \ve} & = \proj_{\va} \mvarM_{\me_2 \comp \ve}
            \end{align*}
        \item $\va' \in \concL{\sigma_2} \ominus \eqz{\concL{\tinstU\me_1}}{\concL{\tinstU\me_2}}$,
              $\va \in \concL{\sigma_3} \ominus \concL{\tinstU\me_2}$, $\va = \concL{\me_1}(\va')$)
              Symmetric to case above.
        \item $\va_1 \in \concL{\sigma_2} \ominus \eqz{\concL{\tinstU\me_1}}{\concL{\tinstU\me_2}}$,
              $\va_2 \in \concL{\sigma_2} \ominus \eqz{\concL{\tinstU\me_1}}{\concL{\tinstU\me_2}}$,
              $\concL{\me_1}(\va_1) = \va = \concL{\me_2}(\va_2)$)
            \begin{align*}
                \proj_\va \mvarM_{\me_1 \comp \eqz{\tinstU\me_1}{\tinstU\me_2}} & = \proj_\va \mvarM_{\me_2 \comp \eqz{\tinstU\me_1}{\tinstU\me_2}} \iff \\
                \sem{\phi(\vC)}\anyL{\me_1}(\proj_{\va_1} \mvarM_{\eqz{\tinstU\me_1}{\tinstU\me_2}}) & = 
                \sem{\phi(\vC)}\anyL{\me_2}(\proj_{\va_2} \mvarM_{\eqz{\tinstU\me_1}{\tinstU\me_2}}) \iff \\
                \sem{\phi(\vC)}\anyL{\me_1}(\proj_{\va_1} \mvarM_\ve) & = 
                \sem{\phi(\vC)}\anyL{\me_2}(\proj_{\va_2} \mvarM_\ve) \iff \\
                \proj_\va \mvarM_{\me_1 \comp \ve} & = \proj_\va \mvarM_{\me_2 \comp \ve}
            \end{align*}
    \end{itemize}


    As a result, by the UP of $\ve$ we have a $\inv^\vh : \eqz{\tinstU\me_1}{\tinstU\me_2} \to \sigma_1$ in $\tinstC$.
    It remains to show that $\inv^\vh$ is the inverse of $\vh$, that is
    \begin{center}
    \begin{tikzcd}
        {\eqz{\tinstU{\me_1}}{\tinstU{\me_2}}} 
            \ar[rr, "{\eqz{\tinstU{\me_1}}{\tinstU{\me_2}}}"]
            \ar[d, bend left=30, dashed, "\inv^\vh" yshift=5pt]
            && {\sigma_2} 
        \\
        {\sigma_1}
            \ar[urr, "\ve"']
            \ar[u, bend left=30, dashed, "\vh" yshift=5pt, "\iso"' yshift=-5pt]
    \end{tikzcd}
    \end{center}
    which follows from $\ve$ and $\eqz{\tinstU\me_1}{\tinstU\me_2}$ being mono.
\end{proof}
% u does not reflect all equalizers necessarily:
% consider \emptyset -{data: True, True}-> 1_{a1,a2 : Bool} -{id, swap}-> 1_{a1,a2 : Bool}
% in T this is an equalizing diagram but in TI we make a choice on the data to be True, True
% which breaks the UP. We could also have given {data: False, False}.

\begin{lemma}[Forgetful Functor Reflects Pullbacks]\label{lem:tinstU-pls-refl}
    $\tinstU : \tinstC \to \tempC$ reflects pullbacks.
\end{lemma}
\begin{proof}
    Consider a diagram 
        \begin{tikzcd}[sep=small]
            & {\sigma} \ar[dl, "\vp_1"'] \ar[dr, "\vp_2"]
        \\
        {\sigma_1} \ar[dr, "\me_1"'] 
                && {\sigma_2} \ar[dl, "\me_2"]
        \\
            & {\hat\sigma}
        \end{tikzcd}
    s.t. its $\tinstU$-image is a pullback.
    Given $\me_1 : \sigma' \to \sigma_1$, $\me_2 : \sigma' \to \sigma_2$ s.t. $\me_1 \comp \me_1' = \me_2 \comp \me_2'$,
    we get a $\vh : \sigma' \to \sigma$ in $\tempC$.

    We equip $\vh$ with data using the canonical iso
    \[ \concL{\sigma_1} \ominus \concL{\me_1'} \iso (\concL{\sigma_1} \ominus \concL{\vp_1}) \oplus (\concL\sigma \ominus \concL\vh) \]
    by giving $\va : \vC \in \concL\sigma \ominus \concL\vh$, the cell $\proj_{\concL{\vp_1}(\va)} \mvarM_{\me_1'}$.
    Notice that we can factor its locations through $\anyL\sigma$ since
    \begin{align*}
        \sem{\phi(\vC)}\anyL{\me_1}(\proj_{\concL{\vp_1}(\va)} \mvarM_{\me_1'}) & = 
        \proj_{\concL{\me_1 \comp \vp_1}(\va)} \mvarM_{\me_1 \comp \me_1'} \\
         & = \proj_{\concL{\me_2 \comp \vp_2}(\va)} \mvarM_{\me_2 \comp \me_2'} \\
         & = \sem{\phi(\vC)}\anyL{\me_2}(\proj_{\concL{\vp_2}(\va)} \mvarM_{\me_2'})
    \end{align*}
    and so for every $\vb : \vC' \in \anyL{\sigma_1}$, we have 
    \[ \vb \isIn \proj_{\concL{\vp_1}(\va)} \mvarM_{\me_1'} \implies \anyL{\me_1}(\vb) = \anyL{\me_2}(\vb) \]
    and so we can factor its locations through $\inv^{\anyL{\vp_1}}$ since its $\tinstU$-image is a pullback. 

    To show that $\me_1' = \vp_1 \comp \vh$, consider $\va : \vC \in \concL{\sigma_1} \ominus \concL{\vp_1}$.
    Since the $\concL\dash$ image is a pullback, it must be that $\concL{\me_1}(\va) \in \concL{\hat\sigma} \ominus \concL{\me_2}$.
    This implies that 
        \[ \sem{\phi(\vC)}\anyL{\me_1}(\proj_\va \mvarM_{\vp_1}) = \proj_\va \mvarM_{\me_1'}
                                                                 = \sem{\phi(\vC)}\anyL{\me_1}(\proj_\va \mvarM_{\me_1'}) \] 
    which implies $\proj_\va \mvarM_{\vp_1} = \proj_\va \mvarM_{\me_1'}$.
    The case for $\me_2' = \vp_1 \comp \vh$ is symmetric and uniqueness of $\vh$ follows from $\proj_1$ being mono.
\end{proof}
% u does not generally preserve pullbacks.
% Consider cospan 1_{a : B} -{a \mapsto a, False}-> 1_{a,b : Bool} <-{a \mapsto a, True} 1_{a : B} 
% In T, the vertex of the pullback is 1_a
% but in TI the data would not match for vertex 1_a.
% So its pullback becomes 1_a <-{a \mapsto a, True}- ?_a -{a \mapsto a, False}-> 1_a
% This is a counterexample that u preserves pullbacks

\begin{lemma}[Location Functor Preserves Pullbacks]\label{lem:anyL-pls-pres}
    $\anyL\dash : \tinstC \to \worldC$ preserves pullbacks.
\end{lemma}
\begin{proof}
    Consider 
    \begin{tikzcd}
        {\sigma} \ar[r, "\vp_2"] \ar[d, "\vp_1"'] \pbsquare{dr}
            & {\sigma_2} \ar[d, "\me_2"]
        \\
        {\sigma_1} \ar[r, "\me_1"']
            & {\hat\sigma}
    \end{tikzcd}.
    If there are $\va_1 : \vC \in \anyL{\sigma_1}$ and $\va_2 : \vC \in \anyL{\sigma_2}$ s.t. $\anyL{\me_1}(\va_1) = \anyL{\me_2}(\va_2)$,
    we can form a new $\sigma' = \concL{\sigma} \to {\anyL\sigma \oplus \va : \vC}$ that maps $\va$ to $\va_1$ and $\va_2$ and notice that the data of $\vp_1$ and $\vp_2$ will be the same.
    As a result, the new square with root $\sigma'$ will commute and thus it must be that $\va : \vC$ is already in $\sigma$.
\end{proof}

\subsubsection{Reachability}

A concept that becomes relevant when we introduce instantiation data with location variables is the notion of reachable and unreachable locations.
We talk about (un)reachable concrete locations and (un)reachable location variables, where reachability refers to whether we can reach the location by traversing the heap starting from a given concrete location.

\begin{definition}[Reachable Locations]\label{def:reach-loc}
    The reachable locations from $\vw \xto\tau \anyL\sigma$ w.r.t. a template heap $\mvarM : \theapF(\sigma)$
    is the least $\vw^\reach \xto{\tau^\ureach} \anyL\sigma$ such that
    \begin{enumerate}[(i)]
        \item World $\vw$ factors through $\vw^\reach$
            \begin{center}
                \begin{tikzcd} {\vw} \ar[dr, "\exists{\tau^\reach}"'] \ar[r, "\tau"] & {\anyL\sigma} \\ 
                                    & {\vw^\reach} \ar[u, "\tau^\ureach"'] 
             k   \end{tikzcd}
            \end{center}
            \item For every $\vb : \vC \in \vw^\reach \pltimes{\anyL\sigma} \concL\sigma$ \ie every reachable concrete location, the locations in cell $\vd = \proj_{\res_{\concL\sigma}\tau^\ureach(\vb)} \mvarM$ are contained in $\vw^\reach$
            \begin{center}
                \begin{tikzcd}
                    {1} \ar[r, "\vd"] \ar[dr, "\exists \vd^\reach"'] &  {\sem{\phi(\vC)}{\anyL\sigma}} \\
                                                                  & {\sem{\phi(\vC)}\vw^\reach} \ar[u, "\sem{\phi(\vC)}{\tau^\ureach}"']
                \end{tikzcd}
            \end{center}
    \end{enumerate}
    Least here means that for any other compatible $\tau' : \vw' \to \anyL\sigma$,
    there exists a unique $\vh : \vw^\reach \to \vw'$ that satisfies $\tau' \comp \vh = \tau^\ureach$.

    We refer to $\vw^\reach$ as the reachable locations and $\anyL\sigma \ominus \tau^\ureach$ as the unreachable ones.
\end{definition}

\begin{example}\label{ex:reach-loc}
    Consider $\tau = \incl^{\va \inclto \va,\vb,\vc}$ and template heap over $1_{\va,\vb,\vc}$ being: 
    \begin{center}
    \renewcommand\ndsize{15}
    \begin{tikzpicture}
        \node[ccell=$\vc$] (c) {};
        \node[ccell=$\va$] (a) [below=0pt of c] {};
        \node[ccell=$\vb$] (b) [below=0pt of a] {};
        \draw[ptr] (c.center) to ++\east to ++\south to (a.east);
        \draw[ptr] ($ (a.center) + (0,-\gap pt) $) to ++\east to ++\south to ($ (b.east) + (0,-\gap pt) $);
    \end{tikzpicture}
    \end{center}
    The least $\tau^\ureach$ is $\incl^{\va,\vb \inclto \va,\vb,\vc}$ since $\vc$ is not reachable from $\va$.
    Notice that $\incl^{\va \inclto \va,\vb,\vc}$ is not the least as it does not satisfy condition (ii) for $\va$.
\end{example}

\paragraph{Reachability in terms of Heap Walks}

We give an alternative, algorithmic description of reachability (\cref{def:reach-loc}) in terms of \emph{heap walks} which describe finite, sequential traversals of the heap.
\begin{definition}[Heap Walks]
    For a heap $\mvarM \in \theapF(\sigma)$, a sequence 
    \[ \va_0 : \vC_0, \hdots, \va_\vn : \vC_\vn, \vc : \vC \in \concL\sigma \times \klstar{\concL\sigma} \times \anyL\sigma \] 
    is called a \emph{heap walk} if
    \[ \All{\vi < \vn} \concL\sigma(\va_{\vi+1}) \isIn \proj_{\va_\vi} \mvarM  \Land \vc \isIn \proj_{\va_\vn} \mvarM \]
    and it contains \emph{no duplicate locations}
    \[ \All{\vi, \vj \le \vn} \vi \neq \vj \implies \va_\vi \neq \va_\vj \Land \All{\vi \le \vn} \concL\sigma(\va_\vi) \neq \vc \]
\end{definition}
We informally use the proposition $(\va \isIn \vd)$ to state that location $\va$ is accessed inside cell $\vd$.
For brevity, we avoid defining this directly but the intuition should be clear.
\begin{example}
    Consider the heap from \cref{ex:reach-loc}. Its heap walks are 
    (i) $\vc, \va, \vb$, 
    (ii) $\vc, \va$, 
    (iii) $\va, \vb$
\end{example}

\begin{definition}[Reachable Location via a Heap Walk]
    A location $\vc : \vC \in \anyL\sigma$ is said to be reachable from $\va : \vC' \in \anyL\sigma$ in heap $\mvarM : \theapF(\sigma)$ if 
    $\va = \vc$ or there exists a heap walk $\va' : \vC', \vec\vb, \vc : \vC$ via $\mvarM$ with $\concL\sigma(\va') = \va$.
\end{definition}

\begin{definition}[Reachable Locations via Heap Walks]\label{def:reach-loc-hwalks}
    The reachable locations from $\vw \xto\tau \anyL\sigma$ w.r.t. a template heap $\mvarM : \theapF(\sigma)$
    is a $\vw^\reach \xto{\tau^\ureach} \anyL\sigma$ that satisfies
    \begin{enumerate}[(i)]
        \item World $\vw$ factors through $\vw^\reach$
            \begin{center}
                \begin{tikzcd} {\vw} \ar[dr, "\exists{\tau^\reach}"'] \ar[r, "\tau"] & {\anyL\sigma} \\ 
                                    & {\vw^\reach} \ar[u, "\tau^\ureach"'] 
                \end{tikzcd}
            \end{center}
        \item For every $\vc : \vC \in \anyL\sigma$, $\vc \in \img{\tau^\ureach} \iff \vc \text{ is reachable from a } \tau(\va)$
        \end{enumerate}
\end{definition}

\begin{lemma}[Reachable Locations Property Equivalence]\label{lem:reach-loc-prop-equiv}
    The two definitions of reachability, \cref{def:reach-loc} and \cref{def:reach-loc-hwalks}, are equivalent.
\end{lemma}
\begin{proof}
    Property (i) of both definitions is identical, so we only show that property (ii) in the two definitions is equivalent.

    \cref{def:reach-loc} $\implies$ \cref{def:reach-loc-hwalks}) 

    Direction $\To$) 
    Let $\vc : \vC \in \vw^\reach$.
    If $\tau^\ureach(\vc) \in \img\tau$ we are done as then $\vc$ is reachable from itself. 
    If $\tau^\ureach(\vc) \not\in \img\tau$, we enumerate all heap walks $\va_0 : \vC_0, \hdots, \va_\vn : \vC_\vn, \tau^\ureach(\vc) : \vC$ of at most $\card{\vw^\reach}$ length, 
    that further satisfy $\All\vi \sigma(\va_\vi) \in \img{\tau^\ureach}$.
    There must exist at least one such heap walk as then $\vw^\reach \setDiff \{\vc\}$ would be the least world.
    If any of these heap walks has $\sigma(\va_0) \in \img\tau$ we are done.  
    Otherwise collect all $\va_\vi$ in all the enumerated heap walks and remove them from $\vw^\reach$, call this $\hat\vw^\reach$.
    We show that $\hat\vw^\reach$ would be the least world which would be a contradiction.
    Take a $\vb \in \hat\vw^\reach \pltimes{\anyL\sigma} \concL\sigma$, and consider the locations in $\vd_\vb$.
    If $\vd_\vb$ contains a location $\vb'$ which is not in $\hat\vw^\reach$ it would mean that $\vb'$ is part of a heap walk generated above. 
    Since the heap walks do not contain duplicates it must be less than $\card{\vw^\reach}$ and thus we can extend that heap walk with $\vb$.
    So that would be a contradiction that $\vb \in \hat\vw^\reach$.
    And so $\hat\vw^\reach$ is the least world which is itself a contradiction and so there must exist at least one heap walk with $\concL\sigma(\va_0) \in \img\tau$.

    Direction $\From$) We have that $\vc : \vC$ is reachable from some $\tau(\va)$. If $\va = \vc$ we are done.
    Otherwise we have a heap walk $\va_0 : \vC_0, \hdots, \va_\vn : \vC_\vn, \vc : \vC$ with $\concL\sigma(\va_0) \in \img\tau$.
    Proceed by induction on $\vn$. If $\vn = 0$, then $\concL\sigma(\va_0) \in \img\tau$, and so it is also in $\img{\tau^\ureach}$.
    If $\va_\vn \in \vw^\reach$ and $\va_{\vn+1} \in \vd_{\va_\vn}$, then $\va_{\vn + 1}$ must also be in $\vw^\reach$ by the  
    fact that $\vd_{\va_\vn}$ factors through $\vw^\reach$.
    So every $\va_\vi$ is in $\vw^\reach$, and thus $\vc$ must also be in $\vw^\reach$.
    \nl
    \cref{def:reach-loc-hwalks} $\implies$ \cref{def:reach-loc}) 
    Consider a $\va : \vC \in \vw^\reach \pltimes{\anyL\sigma} \concL\sigma$, and $\vd_\va \in \sem{\phi(\vC)}\anyL\sigma$.
    Since $\va$ is in $\vw^\reach$, $\va$ is reachable from a $\tau(\va')$ by assumption, and any location in $\vd_\va$ will also be reachable 
    by extending the heap walk with that new location.
    As a result, by the bi-implication, the new location will also be in $\vw^\reach$.
\end{proof}

\begin{lemma}[Decidability of Reachability]\label{lem:reach-loc-dec}
    There exists an algorithm that computes $\vw^\reach \xto{\tau^\ureach} \anyL\sigma$ for every $\vw \xto\tau \anyL\sigma$.
\end{lemma}
\begin{proof} 
    Intuitively, the algorithm will traverse the heap from $\vw \pltimes{\anyL\sigma} \concL\sigma$ and return all the locations found including $\img\tau$. 
\end{proof}

\subsubsection{Connectivity}

We extend the notion of reachability to morphisms in $\tinstC$ to characterise morphisms that satisfy 
(i) all cells in the instantiation data are reachable from locations in its domain, 
or (ii) none of the cells in the instantiation data are reachable from locations in its domain. 
We appropriately call morphisms that satisfy (i) as connected, and (ii) as disconnected.

We first give a factorisation of morphisms in $\tinstC$, by factoring out the reachable and unreachable components of a morphism.

\begin{definition}[Factorisation in $\tinstC$]\label{def:con-fact}
    For a morphism $\me : \sigma_1 \to \sigma_2$, we consider the reachable locations from $\anyL\me$ w.r.t. $\mvarM_\me : \theapF(\concL{\sigma_2} \ominus \concL\me \xto{\compl{\concL\me}} \concL{\sigma_2} \xto{\concL{\sigma_2}} \anyL{\sigma_2})$.
    By \cref{lem:reach-loc-dec} we get 
    \[ \anyL{\sigma_1} \xto{{\anyL\me}^\reach} \vw^\reach \xto{{\anyL\me}^\ureach} \anyL{\sigma_2} \]

    We extend ${\anyL\me}^\reach$ and ${\anyL\me}^\ureach$ into morphisms in $\tinstC$ by taking the \emph{largest possible choice} for the concrete locations,
    that is, we consider the pullback
    \begin{center}
    \begin{tikzcd}
        {\anyL{\sigma_1}} 
            \ar[r, "{\anyL\me}^\reach"] \ar[rrr, bend left=40, "\anyL\me"] 
            & {\vw^\reach} 
                \ar[rr, "{\anyL\me}^\ureach"] 
                && {\anyL{\sigma_2}} 
        \\
        {\concL{\sigma_1}} 
            \ar[u, "\concL{\sigma_1}"] \ar[r, dashed, "\vh"] \ar[rrr, bend right=40, "\concL\me"] 
            & {\vw^\reach \pltimes{\anyL{\sigma_2}} \concL{\sigma_2}} 
                \ar[u] \ar[rr] \pbsquare{urr} 
                && {\concL{\sigma_2}} 
                    \ar[u, "\concL{\sigma_2}"']
    \end{tikzcd}
    \end{center}
    where the instantiation data are given by the canonical isomorphism
    \[ \concL{\sigma_2} \ominus \concL\me \iso (\concL{\sigma_2} \ominus \res_{\concL{\sigma_2}}{\anyL\me}^\ureach) 
                                          \mplus (\vw^\reach \pltimes{\anyL{\sigma_2}} \concL{\sigma_2} \ominus \vh) \]
    
    We show that locations $(\vw^\reach \pltimes{\anyL{\sigma_2}} \concL{\sigma_2}) \ominus \vh$ have cells that access only $\vw^\reach$ from
    \[ (\vw^\reach \pltimes{\anyL{\sigma_2}} \concL{\sigma_2}) \ominus \vh \iso \vw^\reach \pltimes{\anyL{\sigma_2}} (\concL{\sigma_2} \ominus \concL\me) \]
    which follows from pullback pasting
    
    \begin{center}
    \begin{tikzcd}
        {\vw^\reach} 
            \ar[r, "{\anyL\me}^\ureach"] 
            & {\anyL{\sigma_2}} 
        \\
        {\vw^\reach \pltimes{\anyL{\sigma_2}} \concL{\sigma_2}} \pbsquare{ur}
            \ar[u, "\res_{{\anyL\me}^\ureach}\concL{\sigma_2}"]
            \ar[r, "\res_{\concL{\sigma_2}}{\anyL\me}^\ureach"]
            & {\concL{\sigma_2}}
                \ar[u, "\concL{\sigma_2}"']
        \\
        {(\vw^\reach \pltimes{\anyL{\sigma_2}} \concL{\sigma_2}) \ominus \vh} \pbsquare{ur}
            \ar[u, "\compl\vh"]
            \ar[r, "\res_{\concL{\sigma_2}}{\anyL\me}^\ureach"' yshift=-2pt]
            & {\concL{\sigma_2} \ominus \concL\me}
                \ar[u, "\compl{\concL\me}"']
    \end{tikzcd}
    \end{center}

    and so \cref{def:reach-loc} gives us the factoring of the cells through $\vw^\reach$.
\end{definition}

\begin{notation}
    For a morphism $\me : \sigma_1 \to \sigma_2$, we write $\con\me : \sigma_1 \to \sigma_{\con\me}$ and $\uncon\me : \sigma_{\con\me} \to \sigma_2$ for its factorisation given by \cref{def:con-fact}.
    We also write $\sigma_{\uncon\me} \eqdef \sigma_2 \ominus \uncon\me$.
\end{notation}

\begin{example}
    \renewcommand\ndsize{15}

    Consider the morphism $?_\va \xto{\pair{?_{\{\va,\vb\}}}{\incl^{\va \inclto \va,\vb,\vc}}} \incl^{\va,\vb \inclto \va,\vb,\vc}$ with data
    \begin{tikzpicture}[baseline={(a.base)}]
        \node[vcell=$\vc$] (c) {};
        \node[ccell=$\va$] (a) [below=0pt of c] {};
        \node[ccell=$\vb$] (b) [below=0pt of a] {};
        \draw[ptr] (c.center) to ++\east to ++\south to (a.east);
        \draw[ptr] ($ (a.center) + (0,-\gap pt) $) to ++\east to ++\south to ($ (b.east) + (0,-\gap pt) $);
    \end{tikzpicture}

    Its reachable/unreachable factoring given by \cref{def:con-fact} is
    \[ ?_\va \xto{\pair{?_{\{\va,\vb\}}}{\incl^{\va \inclto \va,\vb}}} 1_{\va,\vb} \xto{\pair{\incl^{\va,\vb \inclto \va,\vb,\vc}}{\incl^{\va,\vb \inclto \va,\vb,\vc}}} \incl^{\va,\vb \inclto \va,\vb,\vc} \]
    with data
    \begin{tikzpicture}
        \node[ccell=$\va$] (a) {};
        \node[ccell=$\vb$] (b) [below=0pt of a] {};
        \draw[ptr] ($ (a.center) + (0,-\gap pt) $) to ++\east to ++\south to ($ (b.east) + (0,-\gap pt) $);
    \end{tikzpicture} for the reachable morphism.
\end{example}

\begin{definition}[Connected morphisms]\label{def:connected-morph}
    A morphism $\me : \sigma_1 \to \sigma_2$ is said to be connected if $\uncon\me$ is an isomorphism.
\end{definition}
\begin{definition}[Disconnected morphisms]
    A morphism $\me : \sigma_1 \to \sigma_2$ is said to be disconnected if $\con\me$ is an isomorphism.
\end{definition}

\begin{lemma}[Factoring Connectedness]\label{lem:con-fact-is-con}
    For a morphism $\me : \sigma_1 \to \sigma_2$, its $\con\me$ part is connected.
\end{lemma}
\begin{proof}
    We have the diagram
    \begin{tikzcd}
        {\sigma_1} \ar[rr, "\con\me"] \ar[dr, "\con{\con\me}"'] 
                && {\sigma_{\con\me}} \ar[r, "\uncon\me"] 
                    & {\sigma_2}
        \\
                & {\sigma_{\con{\con\me}}} \ar[ur, "\uncon{\con\me}"']
    \end{tikzcd} and we want to show that $\uncon{\con\me}$ is iso.

    It is enough to show that $\anyL{\uncon{\con\me}}$ is iso, as $\concL{\uncon{\con\me}}$ is given by the pullback
    \begin{center}
    \begin{tikzcd}
        {\anyL{\sigma_{\con{\con\me}}}} 
            \ar[r, "\anyL{\uncon{\con\me}}"] 
            & {\anyL{\sigma_{\con\me}}}
        \\
        {\concL{\sigma_{\con{\con\me}}}} \pbsquare{ur}
            \ar[u] \ar[r, "\concL{\uncon{\con\me}}"']
            & {\concL{\sigma_{\con\me}}} 
                \ar[u, "\concL{\sigma_{\con\me}}"']
    \end{tikzcd}
    \end{center}
    and so if $\anyL{\uncon{\con\me}}$ is iso, so is $\concL{\uncon{\con\me}}$.

    We proceed by showing that 
    \[ \anyL{\sigma_1} \xto{\anyL{\con{\con\me}}} \anyL{\sigma_{\con{\con\me}}} (\xto{\anyL{\uncon{\con\me}}} \anyL{\sigma_{\con\me}} \xto{\anyL{\uncon\me}} \anyL{\sigma_2}) \] 
    satisfies \cref{def:reach-loc} for $\anyL{\sigma_1} \xto{\anyL\me} \anyL{\sigma_2}$ and $\mvarM_\me$.
    As a result, by the least property of \cref{def:reach-loc} for $\anyL{\con\me}$, we get that $\anyL{\uncon{\con\me}}$ must be iso.

    Property (i) is immediate, and so consider property (ii).
    Consider $\va : \vC \in (\concL{\sigma_2} \ominus \concL\me) \pltimes{\anyL{\sigma_2}} \anyL{\sigma_{\con{\con\me}}}$,
    where the pullback is given by
    \begin{center}
        \begin{tikzcd}
        {\anyL{\sigma_{\con{\con\me}}}}
            \ar[r, "\anyL{\uncon{\con\me}}"] 
            & {\anyL{\sigma_{\con\me}}} 
                \ar[r, "\anyL{\uncon\me}"] 
                & {\anyL{\sigma_2}}
        \\
        {\concL{\sigma_{\con{\con\me}}}} \pbsquare{ur}
            \ar[u]
            \ar[r, "\concL{\uncon{\con\me}}"'] 
            & {\concL{\sigma_{\con\me}}} \pbsquare{ur}
                \ar[u]
                \ar[r, "\concL{\uncon\me}"'] 
                & {\concL{\sigma_2}}
                    \ar[u, "\concL{\sigma_2}"']
        \\
        {\concL{\sigma_{\con{\con\me}}} \ominus \concL{\con{\con\me}}} \pbsquare{ur}
            \ar[u, "\compl{\concL{\con{\con\me}}}"]
            \ar[r, "\concL{\uncon{\con\me}}"]
                & {\concL{\sigma_{\con\me}} \ominus \concL{\con\me}} \pbsquare{ur}
                    \ar[u, "\compl{\concL{\con\me}}"]
                    \ar[r, "\concL{\uncon\me}"]
                & {\concL{\sigma_2} \ominus \concL\me}
                    \ar[u, "\compl{\concL\me}"]
        \end{tikzcd}
    \end{center}
    which we use to prove that cell $\va$ has locations in $\anyL{\sigma_{\con\con\me}}$
    \begin{align*}
        \proj_{\concL{\uncon\me}\concL{\uncon{\con\me}}(\va)} \mvarM_\me
            & = \sem{\phi(\vC)}\anyL{\uncon\me}(\proj_{\concL{\uncon{\con\me}}(\va)} \mvarM_{\con\me}) \\
            & = \sem{\phi(\vC)}\anyL{\uncon\me}(\sem{\phi(\vC)}\anyL{\uncon{\con\me}}(\proj_{\va} \mvarM_{\con{\con\me}}))
    \end{align*}
\end{proof}

\begin{lemma}[Factoring and Variable Preservation]\label{lem:con-fact-is-vpres}
    For a morphism $\me : \sigma_1 \to \sigma_2$, its $\uncon\me$ part is variable preserving.
\end{lemma}
\begin{proof}
    Morphism $\uncon\me$ is given by the following diagram
    \begin{center}
    \begin{tikzcd}
        {\anyL{\sigma_1}} \ar[r, "\anyL{\con\me}"] \ar[rr, bend left=40, "\anyL\me"] & {\anyL{\sigma_\me}} \ar[r, "\anyL{\uncon\me}"] & {\anyL{\sigma_2}} \\
        {\concL{\sigma_1}} \ar[u, "\concL{\sigma_1}"] \ar[r, dashed, "\concL{\con\me}"'] \ar[rr, bend right=40, "\concL\me"] & {\concL{\sigma_\me}} \ar[u, "\concL{\sigma_\me}"] \ar[r, "\concL{\uncon\me}"'] \pbsquare{ur} & {\concL{\sigma_2}} \ar[u, "\concL{\sigma_2}"']
    \end{tikzcd}
    \end{center}
    which by \cref{lem:vpres-prop} since $\uncon\me$ is a pullback it is also variable preserving.
\end{proof}

\begin{lemma}[Epimorphisms in $\tinstC$]\label{lem:epis-tinstC}
   A morphism $\me : \sigma_1 \to \sigma_2$ in $\tinstC$ is epi iff $\me$ is connected.
\end{lemma}
\begin{proof}\nl
    $\to$)  
    From \cref{def:con-fact} we can factor $\me$ into $\sigma_1 \xto{\con\me} \sigma_\me \xto{\uncon\me} \sigma_2$ and
    to show that $\me$ is connected it is enough to show that $\anyL{\uncon\me}$ is surjective.

    By \cref{lem:con-fact-is-vpres} we form morphisms $1 \mplus \incl_1$ and $1 \mplus \incl_2$ in $\tinstC$
    \begin{center}
    \begin{tikzcd}
            & {\anyL{\sigma_\me} \mplus (\anyL{\sigma_2} \ominus \anyL{\uncon\me})} 
                \ar[rr, shift left=1, "1 \mplus \incl_1"] \ar[rr, shift right=1, "1 \mplus \incl_2"'] 
                && {\anyL{\sigma_\me} \mplus \Exp{(\anyL{\sigma_2} \ominus \anyL{\uncon\me})}2}
            \\
        {\sigma_2} 
            \ar[r, phantom, "\iso_{\mathsf{r}}" xshift=-28pt] 
            & {}
            \\
            & {\concL{\sigma_\me} \mplus (\concL{\sigma_2} \ominus \concL{\uncon\me})} 
                \ar[uu, "\concL{\sigma_\me} \mplus \concL{\sigma_2}\res|{\uncon\me}"] 
                \ar[rr, shift left=1, "1 \mplus \incl_1"] 
                \ar[rr, shift right=1, "1 \mplus \incl_2"']    
                && {\concL{\sigma_\me} \mplus \Exp{(\concL{\sigma_2} \ominus \concL{\uncon\me})}2}
                    \ar[uu, "\concL{\sigma_\me} \mplus \Exp{\concL{\sigma_2}\res|{\uncon\me}}2"']
    \end{tikzcd}
    \end{center}
    with data $\mvarM_{1 \mplus \incl_\vi} = \dheapF(\iso, 1 \oplus \incl_{\vi \mod 2 + 1} \comp \anyL{\iso_{\mathsf{r}}})(\mvarM_{\uncon\me})$.
    The data $\mvarM_{(1 \oplus \incl_\vi) \comp \uncon\me}$ results in two copies of $\mvarM_{\uncon\me}$ and so
    $\mvarM_{(1 \oplus \incl_1) \comp \uncon\me} = \mvarM_{(1 \oplus \incl_2) \comp \uncon\me}$. 
    As a result 
    \begin{align*}
        (1 \oplus \incl_1) \comp \uncon\me & = (1 \oplus \incl_2) \comp \uncon\me \implies \\
        (1 \oplus \incl_1) \comp \me        & = (1 \oplus \incl_2) \comp \me \implies \\
        1 \oplus \incl_1                    & = 1 \oplus \incl_2 \implies \\
        \anyL{\uncon\me}                   &\text{ is surjective} \implies \\
        \uncon\me                          &\text{ is iso} \\
    \end{align*}

    $\from$) Consider $\me_1 \comp \me = \me_2 \comp \me$.
             We show that $\anyL{\me_1} = \anyL{\me_2}$ which in turn implies $\me_1 = \me_2$.

             Consider $\va' : \vC \in \anyL{\sigma_2}$, which by connectedness gives us a $\va : \vC \in \anyL{\sigma_\me}$ s.t. $\anyL{\uncon\me}(\va) = \va'$.
             If $\va$ is in the image of $\anyL{\con\me}$ we are done.
             Otherwise by \cref{lem:reach-loc-prop-equiv} we have a heap walk $\vb_0 : \vC_0, \hdots, \vb_\vn : \vC_\vn, \va : \vC$
             along $\mvarM_{\sigma_\me}$ where $\concL{\sigma_\me}(\vb_0) \in \img{\anyL{\con\me}}$.

             We proceed by induction on the length of the sequence to show that 
             \[ \All{\vi \le \vn} \anyL{\me_1}\concL{\sigma_\me}(\vb_\vi) = \anyL{\me_2}\concL{\sigma_\me}(\vb_\vi) \]
             $\vn = 0$) Then $\vb_0$ is in the image of $\anyL{\con\me}$ and so $\anyL{\me_1}\concL{\sigma_\me}(\vb_0) = \anyL{\me_2}\concL{\sigma_\me}(\vb_0)$
                        by $\me \comp \me_1 = \me \comp \me_2$.\nl
             $\vn = \vn' + 1$) 
                        We have
                        \begin{align*}
                            \anyL{\me_1}\concL{\sigma_\me}(\vb_{\vn'}) &= \anyL{\me_2}\concL{\sigma_\me}(\vb_{\vn'}) = \tag{$\mathsf{IH}$} \\
                            \concL{\sigma_2}\concL{\me_1}(\vb_{\vn'})  &= \concL{\sigma_2}\concL{\me_2}(\vb_{\vn'}) \implies \tag{by $\concL{\sigma_2}$ is injective} \\
                            \concL{\me_1}(\vb_{\vn'})                  &= \concL{\me_2}(\vb_{\vn'}) \implies        \tag{by $\mvarM_{\me_1 \comp \me} = \mvarM_{\me_2 \comp \me}$} \\
                            \sem{\phi(\vC_{\vn'})}\anyL{\me_1}(\vd_{\vn'})  &= \sem{\phi(\vC_{\vn'})}\anyL{\me_2}(\vd_{\vn'}) \implies  \tag{and since $\vb_\vn \isIn \vd_{\vn'}$} \\
                            \anyL{\me_1}\concL{\sigma_\me}(\vb_\vn)    &= \anyL{\me_2}\concL{\sigma_\me}(\vb_\vn)                       
                        \end{align*}
             Since $\va \text{ is in } \vd_\vn$ and $\anyL{\me_1}\concL{\sigma_\me}(\vb_\vn) = \anyL{\me_2}\concL{\sigma_\me}(\vb_\vn)$,
             we have that $\anyL{\me_1}(\va) = \anyL{\me_2}(\va)$.
\end{proof}

\begin{lemma}[Composition of Connected Morphisms]
    Connected morphisms are closed under composition.
\end{lemma}
\begin{proof}
    Follows immediately from \cref{lem:epis-tinstC} and epis being closed under composition.
\end{proof}

\begin{lemma}[Regular Monos in $\tinstC$]\label{lem:rmono-dcon-tinstC}
    A morphism $\me : \sigma_1 \to \sigma_2$ in $\tinstC$ is regular mono iff $\me$ is variable preserving.
\end{lemma}
\begin{proof}\nl
    $\to$) If $\me$ is part of an equalizer diagram, by \cref{lem:tinstU-eq-pres}, if we forget the data we still have an equalizer diagram in $\tempC$.
           By \cref{lem:rmono-vpres-tempC}, $\tinstU\me$ will be variable preserving and since variable preservation is reflected by $\tinstU$, $\me$ is variable preserving.\nl
    \def\ciso{\iso^{\mathsf{c}}}
    \def\isoO{\iso^1}
    \def\isoT{\iso^2}
    $\from$) 
    We show that this is an equalizer diagram (\cref{lem:vpres-prop} gives us $\concL{\sigma_2}\res|\me$, and $\ciso$ is the canonical isomorphism induced)
    \begin{center}
    \begin{tikzcd}
            & {\anyL{\sigma_1} \mplus (\anyL{\sigma_2} \ominus \anyL\me)} 
                \ar[rr, shift left=1, "1 \mplus \incl_1"] 
                \ar[rr, shift right=1, "1 \mplus \incl_2"'] 
                    && {\anyL{\sigma_1} \mplus (\anyL{\sigma_2} \ominus \anyL\me) \mplus (\anyL{\sigma_2} \ominus \anyL\me)}
\\
        {\sigma_2} \ar[r, phantom, pos=0.20, "\ciso"] & {}
\\
            & {\concL{\sigma_1} \mplus (\concL{\sigma_2} \ominus \concL\me)} 
                \ar[uu, "\concL{\sigma_1} \mplus \concL{\sigma_2}\res|\me"'] 
                \ar[rr, shift left=1, "1 \mplus \incl_1"] 
                \ar[rr, shift right=1, "1 \mplus \incl_2"']    
                    && {\concL{\sigma_1} \mplus (\concL{\sigma_2} \ominus \concL\me) \mplus (\concL{\sigma_2} \ominus \concL\me)} 
                        \ar[uu, "\concL{\sigma_1} \mplus \concL{\sigma_2}\res|\me \mplus \concL{\sigma_2}\res|\me"']
    \end{tikzcd}
    \end{center}
    We equip $1 \oplus \incl_1$ with instantiation data $\dheapF(\isoO, 1 \oplus \incl_2 \comp \anyL\ciso)(\mvarM_\me)$ by noticing the isomorphism
    \[ (\concL{\sigma_1} \mplus (\concL{\sigma_2} \ominus \concL\me) \mplus (\concL{\sigma_2} \ominus \concL\me)) \ominus (1 \oplus \incl_1)
        \isoO \concL{\sigma_2} \ominus \concL\me \]
    We do the same for $1 \oplus \incl_2$, \ie we use $\dheapF(\isoT, 1 \oplus \incl_1 \comp \anyL\ciso)(\mvarM_\me)$.

    To show that $1 \oplus \incl_1 \comp \ciso \comp \me = 1 \oplus \incl_2 \comp \ciso \comp \me$, we need to show that the data match.
    Notice that 
    \[ (\concL{\sigma_1} \mplus \Exp{(\concL{\sigma_2} \ominus \concL\me)}2) \ominus (1 \oplus \incl_1 \comp \concL\me) \iso \Exp{(\concL{\sigma_2} \ominus \concL\me)}2 \]
    and $\mvarM_{1 \oplus \incl_1 \comp \ciso \me} = \mvarM_{1 \oplus \incl_2 \comp \ciso \me}$ as they have two identical copies of $\dheapF(1, \ciso)(\mvarM_\me)$. 

    Now consider $(1 \mplus \incl_1) \comp \ciso \comp \me' = (1 \mplus \incl_2) \comp \ciso \comp \me'$. 
    Then $\tinstU\me'$ must factor through $\tinstU\me$ and $\mvarM_{\me'}$ can be split between the factoring through the canonical isomorphism.
    Uniqueness of this factoring follows from $\me$ being mono.
\end{proof}

\begin{lemma}[Regular Monos are closed under composition in $\tinstC$]\label{lem:rmono-comp}
    Regular monos in $\tinstC$ are closed under composition.
\end{lemma}
\begin{proof}
    Follows from \cref{lem:rmono-dcon-tinstC} and \cref{lem:vpres-dcon-equiv}, $\vu$ preserving and reflecting variable preserving morphisms and \cref{lem:comp-vpres}.
\end{proof}

\begin{lemma}[Variable Preservation and Disconnectedness]\label{lem:vpres-dcon-equiv}
    A morphism $\me : \sigma_1 \to \sigma_2$ is variable preserving iff it is disconnected.
\end{lemma}
\begin{proof}
    The $\from$ direction follows from \cref{lem:con-fact-is-vpres} and since $\con\me$ is iso it is also variable preserving, and variable preservation is closed under composition.
    For the $\to$ direction, by \cref{lem:con-fact-is-vpres} and \cref{lem:vpres-cancl} we have that $\con\me$ is variable preserving.
    Since $\con\me$ is epi and regular mono it is also iso.
\end{proof}

\begin{lemma}[Factoring Disconnectedness]\label{lem:con-fact-is-dcon}
    For a morphism $\me : \sigma_1 \to \sigma_2$, its $\uncon\me$ part is disconnected.
\end{lemma}
\begin{proof}
    By \cref{lem:con-fact-is-vpres}, $\uncon\me$ is variable preserving which by \cref{lem:vpres-dcon-equiv} means it disconnected.
\end{proof}

\subsubsection{Factorisation System}\label{sec:fact-tinstC}

\begin{lemma}[Factorisation is Unique]\label{lem:fact-uniq-tinstC}
    The connectivity (epi-regular mono) factorisation in $\tinstC$ is unique up to unique isomorphism.
\end{lemma}
\begin{proof}
    We show that the $\sigma_1 \xto{\con\me} \sigma_{\con\me} \xto{\uncon\me} \sigma_2$ factorisation is unique up to unique iso, 
    so consider 
    \begin{center}
        \begin{tikzcd}
            {\sigma_1} 
                \ar[r, twoheadrightarrow,  "\con\me"] \ar[rr, twoheadrightarrow, bend left=35, "\me"] \ar[drr, twoheadrightarrow, bend right=70,  "\me_1"'] \ar[dr, "\con{\me_1}"' yshift=2pt]
                & {\sigma_{\con\me}}
                    \ar[-rmonotip, r, "\uncon\me"] \ar[d, dashed, "\vh"', "\iso"]
                    & {\sigma_2}
            \\
                & {\sigma_{\con{\me_1}}}
                    \ar[-rmonotip, r, "\uncon{\me_1}"', "\iso"]
                    & {\sigma'}
                        \ar[-rmonotip, u, "\me_2"']
        \end{tikzcd}
    \end{center}
    and we show that there exists a unique iso $\vh : \sigma_{\con\me} \to \sigma_{\con{\me_1}}$ that commutes in the diagram above.

    We first show that 
    \[ \anyL{\sigma_1} \xto{\anyL{\con{\me_1}}} \anyL{\sigma_{\con{\me_1}}} (\xto{\uncon{\me_1}} \sigma' \xto{\anyL{\me_2}} \sigma_2) \]
    satisfies \cref{def:reach-loc} for $\anyL{\sigma_1} \xto{\anyL\me} \anyL{\sigma_2}$ and $\mvarM_\me$.

    Property (i) is obvious, and for property (ii) consider $(\concL{\sigma_2} \ominus \concL\me) \pltimes{\anyL{\sigma_2}} \anyL{\sigma_{\con{\me_1}}}$, given by
    \begin{center}
        \begin{tikzcd}
            {\anyL{\sigma_{\con{\me_1}}}}
                \ar[r, "\anyL{\uncon{\me_1}}"]
                & {\anyL{\sigma'}}
                    \ar[r, "\anyL{\me_2}"]
                    & {\anyL{\sigma_2}}
            \\
            {\concL{\sigma_{\con{\me_1}}}}
                \pbsquare{ur} \ar[u, "\concL{\sigma_{\con{\me_1}}}"] \ar[r, "\concL{\uncon{\me_1}}"]
                & {\concL{\sigma'}}
                    \pbsquare{ur} \ar[u, "\concL{\sigma'}"] \ar[r, "\concL{\me_2}"]
                    & {\concL{\sigma_2}} 
                        \ar[u, "\concL{\sigma_2}"]
            \\
            {\concL{\sigma_{\con{\me_1}}} \ominus \concL{\con{\me_1}}}
                \pbsquare{ur} \ar[u, "\compl{\concL{\con{\me_1}}}"] \ar[r, "\concL{\uncon{\me_1}}"]
                & {\concL{\sigma'} \ominus \concL{\me_1}}
                    \pbsquare{ur} \ar[u, "\compl{\concL{\me_1}}"] \ar[r, "\concL{\me_2}"]
                & {\concL{\sigma_2} \ominus \concL\me}
                    \ar[u, "\compl{\concL{\me}}"]
        \end{tikzcd}
    \end{center}
    and we show that locations in every $\proj_{\concL{\me_2}\concL{\uncon{\me_1}}(\va)} \mvarM_\me$ factor through $\anyL{\sigma_{\con{\me_1}}}$ 
    \begin{align*}
        \proj_{\concL{\me_2}\concL{\uncon{\me_1}}(\va)} \mvarM_\me & = \proj_{\concL{\me_2}\concL{\uncon{\me_1}}(\va)} \mvarM_{\me_2 \comp \me_1} \\
                                                                    & = \sem{\phi(\vC)}\anyL{\me_2}(\proj_{\concL{\uncon{\me_1}}(\va)} \mvarM_{\me_1}) \\
                                                                    & = \sem{\phi(\vC)}\anyL{\me_2 \comp \uncon{\me_1}}(\proj_\va \mvarM_{\con{\me_1}})
    \end{align*}

    By the UP of reachable locations we have $\anyL\vh : \anyL{\sigma_{\con\me}} \to \anyL{\sigma_{\con{\me_1}}}$, which we extend to $\tinstC$ by
    \begin{center}
    \begin{tikzcd}
        {\anyL{\sigma_{\con\me}}}
            \ar[r, "\anyL\vh"] \ar[rrr, bend left=30, "\anyL{\uncon\me}"]
            & {\anyL{\sigma_{\con{\me_1}}}}
                \ar[r, "\anyL{\uncon{\me_1}}"]
                & {\anyL{\sigma'}}
                    \ar[r, "\anyL{\me_2}"]
                    & {\anyL{\sigma_2}}
        \\
        {\concL{\sigma_{\con\me}}}
            \pbsquare{ur} \ar[r, dashed, "\concL\vh"] \ar[u, "\concL{\sigma_{\con\me}}"] \ar[rrr, bend right=30, "\concL{\uncon\me}"]
            & {\concL{\sigma_{\con{\me_1}}}}
                \pbsquare{ur} \ar[r, "\concL{\uncon{\me_1}}"] \ar[u, "\concL{\sigma_{\con{\me_1}}}"]
                & {\concL{\sigma'}}
                    \pbsquare{ur} \ar[r, "\anyL{\me_2}"] \ar[u, "\concL{\sigma'}"]
                    & {\concL{\sigma_2}}
                        \ar[u, "\concL{\sigma_2}"]
    \end{tikzcd}
    \end{center}
    Notice that since the outer square and the right (combined) square are pullbacks, so is the left square,
    and $\con{\me_1} = \vh \comp \con\me$ and $\uncon\me = \me_2 \comp \uncon{\me_1} \comp \vh$.

    We further show that
    \[ \anyL{\sigma_1} \xto{\anyL{\con\me}} \anyL{\sigma_{\con\me}} (\xto{\anyL\vh} \anyL{\sigma_{\con{\me_1}}} \xto{\anyL{\uncon{\me_2}}} \anyL{\sigma'}) \]
    satisfies \cref{def:reach-loc} for $\anyL{\sigma_1} \xto{\anyL{\me_1}} \anyL{\sigma'}$ and $\mvarM_{\me_1}$.

    Property (i) is obvious, and for property (ii) consider $(\concL{\sigma'} \ominus \concL{\me_1}) \pltimes{\anyL{\sigma'}} \anyL{\sigma_{\con\me}}$, given by 
    \begin{center}
        \begin{tikzcd}
            {\anyL{\sigma_{\con\me}}} \ar[r, "\anyL\vh"] 
                & {\anyL{\sigma_{\con{\me_1}}}}
                    \ar[r, "\anyL{\uncon{\me_2}}"]
                    & {\anyL{\sigma'}}
            \\
            {\concL{\sigma_{\con\me}}}
                \pbsquare{ur} \ar[u, "\concL{\sigma_{\con\me}}"] \ar[r, "\concL\vh"]
                & {\anyL{\sigma_{\con{\me_1}}}}
                    \pbsquare{ur} \ar[u, "\concL{\sigma_{\con{\me_1}}}"] \ar[r, "\concL{\uncon{\me_2}}"]
                    & {\concL{\sigma'}}
                        \ar[u, "\concL{\sigma'}"]
            \\
            {\concL{\sigma_{\con\me}} \ominus \concL{\con\me}}
                \pbsquare{ur} \ar[u, "\compl{\concL{\con\me}}"] \ar[r, "\concL\vh"]
                & {\concL{\sigma_{\con{\me_1}}} \ominus \concL{\con{\me_1}}}
                    \pbsquare{ur} \ar[u, "\compl{\concL{\con{\me_1}}}"] \ar[r, "\concL{\uncon{\me_2}}"]
                    & {\concL{\sigma'} \ominus \concL{\me_1}}
                        \ar[u, "\compl{\concL{\me_1}}"]
        \end{tikzcd}
    \end{center}
    and we show that locations in every $\proj_{\concL{\uncon{\me_1}}\concL\vh(\va)} \mvarM_{\me_1}$ factor through $\anyL{\sigma_{\con\me}}$ 
    \begin{align*}
        \proj_{\concL{\uncon{\me_1}}\concL\vh(\va)} \mvarM_{\me_1} & = \proj_{\concL{\uncon{\me_1}}\concL\vh(\va)} \mvarM_{\uncon{\me_1} \comp \con{\me_1}} \\
                                                                    & = \sem{\phi(\vC)}\anyL{\uncon{\me_1}}(\proj_{\concL\vh(\va)} \mvarM_{\con{\me_1}}) & \tag{def of $\mvarM_{\uncon{\me_1} \comp \con{\me_1}}$}\\
                                                                    & = \sem{\phi(\vC)}\anyL{\uncon{\me_1} \comp \anyL\vh}(\proj_\va \mvarM_{\con\me}) & \tag{UP of $\anyL{\sigma_{\con\me}}$}
    \end{align*}

    As a result, we get a $\anyL{\inv^\vh} : \anyL{\sigma_{\con{\me_1}}} \to \anyL{\sigma_{\con\me}}$ which is the inverse of $\anyL\vh$ by
    \begin{align*}
        \anyL\vh \comp \anyL{\inv^\vh} & = 1 \iff \\
        \anyL{\uncon\me_1} \comp \anyL\vh \comp \anyL{\inv^\vh} & = \anyL{\uncon{\me_1}} \iff \\
        \anyL{\me_2} \comp \anyL{\uncon\me_1} \comp \anyL\vh \comp \anyL{\inv^\vh} & = \anyL{\me_2} \comp \anyL{\uncon{\me_1}} \iff \\
        \anyL{\uncon\me} \comp \anyL{\inv^\vh} & = \anyL{\me_2} \comp \anyL{\uncon{\me_1}}
    \end{align*}
    and $\anyL\vh \comp \anyL{\inv^\vh} \comp \anyL\vh = \anyL\vh \iff \anyL{\inv^\vh} \comp \anyL\vh = 1$.

    We extend $\anyL{\inv^\vh}$ into a morphism in $\tempC$ by
    \begin{center}
    \begin{tikzcd}
        {\anyL{\sigma_{\con{\me_1}}}} 
            \ar[r, "\anyL{\inv^\vh}"] \ar[rr, bend left=30, "\anyL{\me_2} \comp \anyL{\uncon{\me_1}}"]
            & {\anyL{\sigma_{\con\me}}}
                \ar[r, "\anyL{\uncon\me}"]
                    & {\anyL{\sigma_2}}
        \\
        {\concL{\sigma_{\con{\me_1}}}} 
            \pbsquare{ur} \ar[u, "\concL{\sigma_{\con{\me_1}}}"] \ar[r, dashed, "\concL{\inv^\vh}"] \ar[rr, bend right=30, "\concL{\me_2} \comp \concL{\uncon{\me_1}}"']
            & {\concL{\sigma_{\con\me}}}
                \pbsquare{ur} \ar[u, "\concL{\sigma_{\con\me}}"] \ar[r, "\concL{\uncon\me}"]
                    & {\concL{\sigma_2}}
                        \ar[u, "\concL{\sigma_2}"]
    \end{tikzcd}
    \end{center}
    where the left square is a pullback because the outer and right square are pullbacks, and regular monos compose (\cref{lem:rmono-comp}).
    As a result, $\concL{\inv^\vh}$ is also iso.
    
    Uniqueness of iso $\vh$ follows from uniqueness of the underlying components.  
\end{proof}

\begin{lemma}[Orthogonal Factorisation System in $\tinstC$]\label{lem:ofs-tinstC}
    The classes of epis (connected) and regular monos (disconnected) morphisms gives rise to an orthogonal factorisation system in $\tinstC$.
\end{lemma}
\begin{proof}
    Follows from \cref{lem:fact-uniq-tinstC}, 
    epis being closed under composition and including all isomorphisms,
    and regular monos in $\tinstC$ being closed under composition (\cref{lem:rmono-comp}) and include all isomorphisms.
\end{proof}
 
\subsubsection{Unification}
\label{sec:unification-theaps}

A natural question to ask when dealing with two template heaps that have a common set of public locations is whether they unify,
that is, if there exist instantiations that we can apply to both heaps respectively such that they result in the same heap.
Through the isomorphism $\theapF \iso \homF\tinstC{?_\emptyset}\dash$, we simply phrase this question as whether  
a span $\sigma_1 \xfrom{\me_1} \sigma \xto{\me_2} \sigma_2$ has a cospan $\sigma_1 \xto{\me_1'} \hat\sigma \xfrom{\me_2'} \sigma_2$ such that
the obvious square commutes.

\begin{definition}[Unifier of Morphisms in $\tinstC$]
    We say that two morphisms $\sigma_1 \xfrom{\me_1} \sigma \xto{\me_2} \sigma_2$ unify if there exists a diagram
    \begin{center}
    \begin{tikzcd}
                                      & {\sigma} \ar[dl, "\me_1"'] \ar[dr, "\me_2"] & \\
        {\sigma_1} \ar[dr, "\me_1'"'] &                                             & {\sigma_2} \ar[dl, "\me_2'"] \\
                                      & {\hat\sigma}                                &
    \end{tikzcd}
    \end{center}
    We call $\me_1', \me_2'$ a unifier of $\me_1, \me_2$.
\end{definition}
Notice that some spans do not have any compatible cospan, that is, some pair of heaps do not unify as the following example shows.
\begin{example}
    Consider the following diagram where both morphisms instantiate the variable $\va$ but with different values (true (T) and false (F)).
    \begin{center}
    \begin{tikzcd}[remember picture]
        & {?_\va} \ar[dl] \ar[dr] & \\ 
        |[alias=s1]|{1_\va} & & 
        |[alias=s2]|{1_\va} & & 
    \end{tikzcd}
    \renewcommand\ndsize{9}
    \begin{tikzpicture}[overlay, remember picture] 
        \node[ccell=$\va$] (d1) [above=15pt of s1, xshift=15pt] {\tiny T};
        \node[ccell=$\va$] (d2) [above=15pt of s2, xshift=-15pt] {\tiny F};
    \end{tikzpicture}
    \end{center}
    There cannot be any unifiers that would make the diagram commute as the resulting heaps will always be different.
\end{example}

\paragraph{General Unifiers}

We seek to identify the most general unifers of two morphisms $\sigma_1 \xfrom{\me_1} \sigma \xto{\me_2} \sigma_2$,
that is, the least cospan $\sigma_1 \xto{\me_1'} \hat\sigma \xfrom{\me_2'} \sigma_2$ such that every other cospan factors uniquely thorough it. 
Unfortunately such M.G.U. does not always exist as the following counterexample shows.
\begin{example}\label{ex:mgu-counterexample}
    Consider the following two unifers\nl\nl

    \renewcommand\ndsize{9}
    \hspace{-10pt}
    \begin{tikzcd}[remember picture]
        & {?_\va} \ar[dl] \ar[dr] & \\ 
        |[alias=s1]|{\incl^{\va,\vb,\vd \inclto \va,\vb,\vc,\vd}} \ar[dr] & & 
        |[alias=s2]|{\incl^{\va,\vc,\vd \inclto \va,\vb,\vc,\vd}} \ar[dl] & & \\
        & |[alias=hs]|{1_{\va,\vb,\vc,\vd}} 
    \end{tikzcd}\hspace{-50pt}
    \begin{tikzcd}[remember picture]
        & {?_\va} \ar[dl] \ar[dr] & \\ 
        |[alias=t1]|{\incl^{\va,\vb,\vd \inclto \va,\vb,\vc,\vd}} \ar[dr] & & 
        |[alias=t2]|{\incl^{\va,\vc,\vd \inclto \va,\vb,\vc,\vd}} \ar[dl] & & \\
        & |[alias=ht]|{1_{\va,\vb,\vc,\vd_1, \vd_2}}
    \end{tikzcd}

    \begin{tikzpicture}[overlay, remember picture] 
        \node[ccell=$\vd$] (sd1) [above=35pt of s1, xshift=30pt] {};
        \node[ccell=$\vb$] (sb1) [below=0pt of sd1] {};
        \node[ccell=$\va$] (sa1) [below=0pt of sb1] {};
        \node[vcell=$\vc$] (sc1) [below=0pt of sa1] {};
        \draw[ptr] (sb1.center) to ++\east to ++\north to (sd1.east);
        \draw[ptr] ($ (sa1.center) + (0,\gap pt) $) to ++\east to ++\nnnorth to ($ (sb1.east) + (0,-\gap pt) $);
        \draw[ptr] ($ (sa1.center) + (0,-\gap pt) $) to ++\east to ++\ssouth to (sc1.east);
    
        \node[vcell=$\vb$] (sb2) [above=35pt of s2, xshift=-30pt] {};
        \node[ccell=$\va$] (sa2) [below=0pt of sb2] {};
        \node[ccell=$\vc$] (sc2) [below=0pt of sa2] {};
        \node[ccell=$\vd$] (sd2) [below=0pt of sc2] {};
        \draw[ptr] ($ (sa2.center) + (0,\gap pt) $) to ++\east to ++\nnorth to (sb2.east);
        \draw[ptr] ($ (sa2.center) + (0,-\gap pt) $) to ++\east to ++\sssouth to ($ (sc2.east) + (0,\gap pt) $);
        \draw[ptr] (sc2.center) to ++\east to ++\south to (sd2.east);

        \node[ccell=$\vc$] (sc1p) [left=15pt of hs,yshift=2pt] {};
        \draw[ptr] (sc1p.center) to ++\east to node[midway,xshift=\gap pt, scale=\shpscale] {$\vd$} ++\north;

        \node[ccell=$\vb$] (sb2p) [right=15pt of hs,yshift=5pt] {};
        \draw[ptr] (sb2p.center) to ++\east to node[midway,xshift=\gap pt, scale=\shpscale] {$\vd$} ++\south;

        \node[ccell=$\vd$] (td1) [above=35pt of t1, xshift=30pt] {};
        \node[ccell=$\vb$] (tb1) [below=0pt of td1] {};
        \node[ccell=$\va$] (ta1) [below=0pt of tb1] {};
        \node[vcell=$\vc$] (tc1) [below=0pt of ta1] {};
        \draw[ptr] (tb1.center) to ++\east to ++\north to (td1.east);
        \draw[ptr] ($ (ta1.center) + (0,\gap pt) $) to ++\east to ++\nnnorth to ($ (tb1.east) + (0,-\gap pt) $);
        \draw[ptr] ($ (ta1.center) + (0,-\gap pt) $) to ++\east to ++\ssouth to (tc1.east);
    
        \node[vcell=$\vb$] (tb2) [above=35pt of t2, xshift=-30pt] {};
        \node[ccell=$\va$] (ta2) [below=0pt of tb2] {};
        \node[ccell=$\vc$] (tc2) [below=0pt of ta2] {};
        \node[ccell=$\vd$] (td2) [below=0pt of tc2] {};
        \draw[ptr] ($ (ta2.center) + (0,\gap pt) $) to ++\east to ++\nnorth to (tb2.east);
        \draw[ptr] ($ (ta2.center) + (0,-\gap pt) $) to ++\east to ++\sssouth to ($ (tc2.east) + (0,\gap pt) $);
        \draw[ptr] (tc2.center) to ++\east to ++\south to (td2.east);

        \node[ccell=$\vc$] (tc1p) [left=15pt of ht,yshift=5pt] {};
        \node[ccell=$\vd_2$] (td12p) [below=0pt of tc1p] {};
        \draw[ptr] (tc1p.center) to ++\east to ++\south to (td12p.east);

        \node[ccell=$\vd_1$] (td21p) [right=15pt of ht,yshift=5pt] {};
        \node[ccell=$\vb$] (tb2p) [below=0pt of td21p] {};
        \draw[ptr] (tb2p.center) to ++\east to ++\north to (td21p.east);
    \end{tikzpicture}

    \renewcommand\ndsize{15}
    The unifier on the left results in the heap
    \begin{tikzpicture}[baseline={(b1.base)}]
        \node[ccell=$\vd$] (d1) [] {};
        \node[ccell=$\vb$] (b1) [below=0pt of d1] {};
        \node[ccell=$\va$] (a1) [below=0pt of b1] {};
        \node[ccell=$\vc$] (c1) [below=0pt of a1] {};
        \draw[ptr] (b1.center) to ++\east to ++\north to (d1.east);
        \draw[ptr] ($ (a1.center) + (0,\gap pt) $) to ++\east to ++\nnnorth to ($ (b1.east) + (0,-\gap pt) $);
        \draw[ptr] ($ (a1.center) + (0,-\gap pt) $) to ++\east to ++\sssouth to ($ (c1.east) + (0,\gap pt) $);
        \draw[ptr] (c1.center) to ++\eeast to ++\north to ++\north to ++\north to ++(0,\gap pt) to ($ (d1.east) + (0,\gap pt) $);
    \end{tikzpicture}
    whereas the one on the right results in
    \begin{tikzpicture}[baseline={(b1.base)}]
        \node[ccell=$\vd_1$] (d1) [] {};
        \node[ccell=$\vb$] (b1) [below=0pt of d1] {};
        \node[ccell=$\va$] (a1) [below=0pt of b1] {};
        \node[ccell=$\vc$] (c1) [below=0pt of a1] {};
        \node[ccell=$\vd_2$] (d2) [below=0pt of c1] {};
        \draw[ptr] (b1.center) to ++\east to ++\north to (d1.east);
        \draw[ptr] ($ (a1.center) + (0,\gap pt) $) to ++\east to ++\nnnorth to ($ (b1.east) + (0,-\gap pt) $);
        \draw[ptr] ($ (a1.center) + (0,-\gap pt) $) to ++\east to ++\sssouth to ($ (c1.east) + (0,\gap pt) $);
        \draw[ptr] ($ (c1.center) + (0,-\gap pt) $) to ++\east to ++\sssouth to ($ (d2.east) + (0,\gap pt) $);
    \end{tikzpicture}
    which are mutually incomparable.
\end{example}

Even though most general unifiers do not always exist, we do have that every unifier can be factored though "special" unifiers that are themselves mutually incomparable. 
We call these special unifiers \emph{general unifiers}.

\begin{definition}[General Unifiers]\label{def:general-unifiers}
    We call a unifier $\sigma_1 \xto{\me_1'} \hat\sigma \xfrom{\me_2'} \sigma_2$ general if 
    $\concL{\sigma_1} \xto{\concL{\me_1'}} \concL{\hat\sigma} \xfrom{\concL{\me_2'}} \concL{\sigma_2}$ are jointly surjective and
    $\anyL{\sigma_1} \xto{\anyL{\me_1'}} \anyL{\hat\sigma} \xfrom{\anyL{\me_2'}} \anyL{\sigma_2}$ are jointly surjective.
\end{definition}
\begin{remark}
    By \cref{lem:gen-cospan-epi-prop}, this is equivalent to saying that a unifier is general if it is generated by a jointly surjective cospan
    $\anyL{\sigma_1} \xto{\anyL{\me_1'}} \anyL{\hat\sigma} \xfrom{\anyL{\me_2'}} \anyL{\sigma_2}$.
\end{remark}

\begin{example}
    Both unifiers from \cref{ex:mgu-counterexample} are general unifiers.
    The unifier on the left is generated by 
        \[ \va,\vb,\vc,\vd \xto{1} \va,\vb,\vc,\vd \xfrom{1} \va,\vb,\vc,\vd \]
    whereas the one on the right by 
        \[ \va,\vb,\vc \mplus \vd \xto{1 \mplus \incl_1} \va,\vb,\vc,\vd_1,\vd_2 \xfrom{1 \mplus \incl_2} \va,\vb,\vc \mplus \vd \]
\end{example}

\begin{lemma}[Data of General Unifiers]\label{lem:general-unifier-data}
    If 
    \begin{tikzcd}[sep=small]
        & {\sigma} 
            \ar[dl, "\me_1"'] \ar[dr, "\me_2"] 
    \\ 
    {\sigma_1} 
        \ar[dr, "\me_1'"']
            && {\sigma_2} 
                \ar[dl, "\me_2'"] 
    \\ 
        & {\hat\sigma}
    \end{tikzcd}
    is a general unifier then $\mvarM_{\me_1'}$ and $\mvarM_{\me_2'}$
    are completely determined by $\mvarM_{\me_1}$ and $\mvarM_{\me_2}$.
\end{lemma}
\begin{proof}
    W.l.g. we show that $\mvarM_{\me_2'}$ is completely determined by $\mvarM_{\me_1}$ and the other direction is symmetric.
    Since $\concL{\me_1'}$ and $\concL{\me_2'}$ are jointly surjective, every $\va : \vC \in \concL{\hat\sigma} \ominus \concL{\me_2'}$
    satisfies that there exists a unique $\va' : \vC \in \concL{\sigma_1} \ominus \concL{\me_1}$ s.t. $\concL{\me_1'}(\va') = \va$. 

    This implies by definition of $\mvarM_{\me_2' \comp \me_2}$ and $\mvarM_{\me_2' \comp \me_2} = \mvarM_{\me_1' \comp \me_1}$ that
    \[ \proj_\va \mvarM_{\me_2'} = \proj_\va \mvarM_{\me_2' \comp \me_2} = \proj_{\concL{\me_1'}(\va')} \mvarM_{\me_1' \comp \me_1} = \sem{\phi(\vC)}\anyL{\me_1'}(\proj_{\va'} \mvarM_{\me_1}) \]
\end{proof}

\begin{lemma}[General Unifier Property]\label{lem:general-unifier-prop}
    Every unifier of $\sigma_1 \xfrom{\me_1} \sigma \xto{\me_2} \sigma_2$ factors uniquely (up to iso) through a general unifier.
\end{lemma}
\begin{proof}
    For a unifier $\sigma_1 \xto{\me_1'} \sigma' \xfrom{\me_2'} \sigma_2$ we show the following diagram:
    \begin{center}
    \begin{tikzcd}
            & \sigma \ar[dl, "\me_1"'] \ar[dr, "\me_2"] & 
        \\
        {\sigma_1} \ar[ddr, "\me_1'"'] \ar[dr, "\hat\me_1"]
            && {\sigma_2} \ar[ddl, "\me_2'"] \ar[dl, "\hat\me_2"']
        \\
            & {\hat\sigma} \ar[d, dashed, "\vh" yshift=5pt]
        \\
            & {\sigma'}
    \end{tikzcd}
    \end{center}
    We define the general unifier $\sigma_1 \xto{\hat\me_1} \hat\sigma \xfrom{\hat\me_2} \sigma_2$ as follows.
    \begin{enumerate}
        \item We define $\anyL{\hat\me_1}$ and $\anyL{\hat\me_2}$ as
              \begin{center}
              \begin{tikzcd}
                  {\anyL{\sigma_1}} \ar[dr, dashed, "\anyL{\hat\me_1}"'] \ar[r, "\anyL{\me_1'}"] 
                      & {\anyL{\sigma'}} 
                          & {\anyL{\sigma_2}} 
                              \ar[dl, dashed, "\anyL{\hat\me_2}"] \ar[l, "\anyL{\me_2'}"'] \\ 
                          & {\anyL{\hat\sigma} \eqdef \img{\anyL{\me_1'}} \setUnion \img{\anyL{\me_2'}}} \ar[u, hook, "\anyL\vh"']
              \end{tikzcd}
              \end{center}
              It is clear that they are jointly surjective.
        \item We generate morphisms $\tinstU\hat\me_1$ and $\tinstU\hat\me_2$ from $\anyL{\hat\me_1}$ and $\anyL{\hat\me_2}$ using \cref{lem:gen-cospan},
              and so by the universal property of generated morphisms we have the following diagram in $\tempC$
              \begin{center}
              \begin{tikzcd}
                      & \sigma \ar[dl, "\tinstU\me_1"'] \ar[dr, "\tinstU\me_2"] & 
                  \\
                  {\sigma_1} \ar[ddr, "\tinstU\me_1'"'] \ar[dr, "\tinstU\hat\me_1"]
                      && {\sigma_2} \ar[ddl, "\tinstU\me_2'"] \ar[dl, "\tinstU\hat\me_2"']
                  \\
                      & {\hat\sigma} \ar[d, dashed, "\tinstU\vh" yshift=5pt]
                  \\
                      & {\sigma'}
              \end{tikzcd}
              \end{center}
              That $\tinstU\hat\me_1 \comp \tinstU\me_1 = \tinstU\hat\me_2 \comp \tinstU\me_2$ follows from the outer square and inner triangles commuting.
        \item For the instantiation data we show that $\mvarM_{\me_\vi'}$ can be factored into $\mvarM_\vh$ and $\mvarM_{\hat\me_\vi}$. 
              W.l.g. let $\vi = 1$.
              Through the canonical isomorphism
              \[ \concL{\sigma'} \ominus \concL{\vu\me_1'} \iso (\concL{\sigma'} \ominus \concL{\vu\vh}) \mplus (\concL{\hat\sigma} \ominus \concL{\vu\hat\me_1}) \]
              we have heaps
\[ \mvarM_{\hat\me_1} : \dheapF(\concL{\hat\sigma} \ominus \concL{\hat\me_1}, \anyL{\sigma'}) \;\text{ and }\; \mvarM_\vh : \dheapF(\concL{\sigma'} \ominus \concL\vh, \anyL{\sigma'}) \]
              We further show for every $\va : \vC \in \concL{\hat\sigma} \ominus \hat\me_1$ that $\vd = \proj_\va \mvarM_{\hat\me_1}$ satisfies
              \begin{center}
                  \begin{tikzcd}
                      {1} \ar[dr, "\exists \hat\vd"'] \ar[r, "\vd"] & {\sem{\phi(\vC)}\anyL{\sigma'}} \\
                                                                    &  {\sem{\phi(\vC)}\anyL{\hat\sigma}} \ar[u, "\anyL\vh"']
                  \end{tikzcd}
              \end{center}
              By \cref{lem:gen-cospan-epi-prop} we have that $\concL{\hat\me_1}$ and $\concL{\hat\me_2}$ are jointly surjective, 
              and since $\va : \vC \in \concL{\hat\sigma} \ominus \concL{\hat\me_1}$, and $\vu\hat\me_1 \comp \vu\me_1 = \vu\hat\me_2 \comp \vu\me_2$ we must have a $\va' : \vC \in \concL{\sigma_2} \ominus \concL{\me_2}$ s.t. $\concL{\hat\me_2}(\va') = \va$.
              And so
              \begin{align*}
                  \proj_\va \mvarM_{\hat\me_1} &= \proj_{\concL\vh(\va)} \mvarM_{\me_1'} \\
                                               &= \proj_{\concL\vh(\va)} \mvarM_{\me_1' \comp \me_1} \\
                                               &= \proj_{\concL\vh(\va)} \mvarM_{\me_2' \comp \me_2} \\
                                               &= \proj_{\concL\vh(\concL{\hat\me_2}(\va'))} \mvarM_{\me_2' \comp \me_2} \\
                                               &= \proj_{\concL{\me_2'}(\va')} \mvarM_{\me_2' \comp \me_2} \\
                                               &= \sem{\phi(\vC)}\anyL{\me_2'}(\proj_{\va'} \mvarM_{\me_2}) \\
                                               &= \sem{\phi(\vC)}\anyL\vh(\sem{\phi(\vC)}\anyL{\hat\me_2}(\proj_{\va'} \mvarM_{\me_2}))
              \end{align*}

              Lastly, $\mvarM_{\hat\me_1 \comp \me_1} = \mvarM_{\hat\me_2 \comp \me_2}$ follows from $\vh$ being mono (\cref{lem:monos-tinistC}).
    \end{enumerate}
    To show uniqueness, consider another general unifier $\sigma_1 \xto{\check\me_1} \check\sigma \xfrom{\check\me_2} \sigma_2$.

    Notice that the induced morphism from $\check\sigma \to \sigma'$ must factor through $\hat\sigma$ in an isomorphic way
    \begin{center}
        \begin{tikzcd}
            {\sigma_1}
                \ar[r, "\hat\me_1"]
                    \ar[dr, "\check\me_1"']
                    & {\hat\sigma}
                        & {\sigma_2}
                            \ar[l, "\hat\me_2"']
                                \ar[dl, "\check\me_2"]
            \\
                    & {\check\sigma}
                        \ar[u, dashed, "\iso"]
        \end{tikzcd}
    \end{center}
    It remains to show that their data match which follows from \cref{lem:general-unifier-data}.
\end{proof}

\begin{lemma}[General Unifiers of Identities]\label{lem:general-unifier-id}
    The general unifiers of a span $\sigma \xfrom{1} \sigma \xto\me \sigma'$ is only cospan $\sigma \xto\me \sigma' \xfrom{1} \sigma'$ (up to iso).
\end{lemma}
\begin{proof}
    To show that cospan $\sigma \xto\me \sigma' \xfrom{1} \sigma'$ is a general unifier, we need it to be jointly surjective at the level $\concL\dash$ and $\anyL\dash$, which is immediate from the identity morphism.
    That every other unifier factors through this (and so it is unique up to iso), is immediate.
\end{proof}

\begin{lemma}[Pasting Law for General Unifiers]\label{lem:general-unifier-pasting-law}
    In a diagram
    \begin{tikzcd}
        {\sigma} \ar[r, "\me_1"] \ar[d, "\me_1'"']
            & {\sigma_1} \ar[r, "\me_3"] \ar[d, "\me_2"]
                & {\sigma_3} \ar[d, "\me_3'"]
        \\
        {\sigma_1'} \ar[r, "\me_2'"']
            & {\sigma_2} \ar[r, "\me_3''"']
                & {\sigma_3'}
    \end{tikzcd}
    where the left square is a general unifier,  
    then the whole square is a general unifier iff the right square is.
\end{lemma}
\begin{proof}
    For the $\to$ direction, it is clear that if any $\anyL{\me_3''}$ and $\anyL{\me_3'}$ are not jointly surjective then the whole square cannot be as well which is a contradiction.
    The same applies to the $\concL\dash$ level, and so the right square must be a general unifier.
    \nl
    For the $\from$ direction, consider a $\va : \vC \in \anyL{\sigma_3'}$.
    Then by the right square being a general unifier, either $\va$ is mapped from $\anyL{\sigma_3}$ or it comes from $\anyL{\sigma_2}$.
    If it comes from $\anyL{\sigma_2}$, then by the left square being a general unifier either it is mapped from $\sigma_1'$ or it is mapped from $\sigma_1$ and since it is a commuting square, it must also be mapped from $\sigma_3$.
    The same reasoning applies to $\concL\dash$.
\end{proof}

\begin{lemma}[General Unifiers Preserve Connectivity]\label{lem:general-unifier-con-prop}
    If 
    \begin{tikzcd}[sep=small]
        & {\sigma} 
            \ar[dl, twoheadrightarrow, "\me_1"'] \ar[dr, "\me_2"] 
    \\ 
    {\sigma_1} 
        \ar[dr, "\me_1'"']
            && {\sigma_2} 
                \ar[dl, "\me_2'"] 
    \\ 
        & {\hat\sigma}
    \end{tikzcd}
    is a general unifier and $\me_1$ is connected, then so is $\me_2'$.
\end{lemma}
\begin{proof}
    Consider a $\va : \vC \in \anyL{\hat\sigma} \ominus \anyL{\uncon{\me_2'}}$.
    This means that $\va : \vC \in \anyL{\hat\sigma} \ominus \anyL{\me_2'}$, and so there exists $\va' : \vC \in \anyL{\sigma_1}$ s.t.
    $\anyL{\me_1'}(\va') = \va$ by definition of general unifiers (\cref{def:general-unifiers}).
    From the factorisation we have a diagram which leads to a contradiction that it commutes unless $\anyL{\hat\sigma} \ominus \anyL{\uncon{\me_2'}} = \emptyset$.
    \begin{center}
        \begin{tikzcd}
            & {\sigma} 
                \ar[dl, twoheadrightarrow, "\me_1"'] \ar[d, twoheadrightarrow, "\con{\me_2' \comp \me_2}" yshift=-3pt] \ar[dr, "\me_2"]
        \\
        {\sigma_1}
            \ar[ddr, "\me_1'"'] \ar[r, twoheadrightarrow, dashed]
            & {\sigma_{\con{\me_2' \comp \me_2}}}
                \ar[dd, -rmonotip] \ar[rd, -rmonotip, dashed]
                & {\sigma_2}
                    \ar[d, twoheadrightarrow, "\con{\me_2'}"]
        \\
                && {\sigma_{\con{\me_2'}}} 
                    \ar[dl, -rmonotip, "\uncon{\me_2'}"]
        \\
            & {\hat\sigma}
        \end{tikzcd}
    \end{center}
    Since $\anyL{\uncon{\me_2'}}$ is iso and $\uncon{\me_2'}$ is a regular mono, we have that $\concL{\uncon{\me_2'}}$ is iso and thus connected.
\end{proof}


\begin{lemma}[Decidability of Unification]\label{lem:unification-dec}
    There exists an algorithm that decides for a span $\sigma_1 \xfrom{\me_1} \sigma \xto{\me_2} \sigma_2$ if they unify or not.
    If they unify, it returns a non-empty finite set of general unifiers
    $\GU{\me_1}{\me_2} \subseteqfin \Coprod{\hat\sigma \in \objC\tempC}{\homF\tinstC{\sigma_1}{\hat\sigma} \times \homF\tinstC{\sigma_2}{\hat\sigma}}$
    that contains all its general unifiers up to isomorphism.
\end{lemma}
\begin{proof}
    By \cref{lem:general-unifier-prop} we have that every unifier factors through a general one and there are only a finite number of isomorphic classes of general unifiers by definition of their jointly surjective property.
    As a result, the algorithm will enumerate every isomorphic class of jointly surjective cospans $\anyL{\sigma_1} \xto{\tau_1} \hat\vw \xfrom{\tau_2} \anyL{\sigma_2}$
    by considering a representative for each class (note that $\objC\worldC$ is countable when $\typeSort$ is, which makes it possible to constructively pick one). 

    It will then equip every representative cospan with data given by \cref{lem:gen-cospan-data} for every side which by \cref{lem:general-unifier-data} tells us that it should be the only data that we have to consider. 
    All that is left to check for it to be a general unifier is that it forms a commuting square in $\tinstC$.
    If the algorithm finds no such commuting generated cospans then it must be that the span does not unify.
\end{proof}

\paragraph{Unification and Local Independence Structure}\nl

\begin{definition}[Local Independent Unifiers]
    A unifier is called a local independent unifier (L.I.U.) if its $\tinstU$ image is a local independent square.
\end{definition}

\begin{definition}[Local Independent General Unifiers]
    A unifier is called a local independent general unifier (L.I.G.U.) if its $\tinstU$ image is a local independent coproduct square.
\end{definition}

\begin{lemma}[Local Independent General Unifier is General]
    A local independent general unifier is a general unifier.
\end{lemma}
\begin{proof}
    To show that a L.I.G.U. $\sigma_1 \xto{\me_1'} \hat\sigma \xto{\me_2'} \sigma_2$ 
    is a general unifier it is enough to show that it is generated by
    $\anyL{\sigma_1} \xto{\anyL{\me_1'}} \anyL{\hat\sigma} \xfrom{\anyL{\me_2'}} \anyL{\sigma_2}$ 
    and moreover this cospan is jointly surjective.
    This is the case as L.I.C. in $\tempC$ is generated from the L.I.C. in $\worldC$ which is itself a jointly surjective cospan. 
\end{proof}

\begin{lemma}[Uniqueness of Local Independent General Unifiers]\label{lem:ligu-uniq}
    L.I.G.U. are unique up to iso.
\end{lemma}
\begin{proof}
    Any two L.I.G.U. are iso in $\tempC$, by L.I.C. being unique up to iso.
    That their data must also be match follows from \cref{lem:general-unifier-data}.
\end{proof}

\begin{lemma}[Local Independent General Unifiers Property]\label{lem:ligu-prop}
    The general unifier of a local independent unifier is a local independent general unifier.
\end{lemma}
\begin{proof}
    The general unifier of a L.I.U. $\sigma_1 \xto{\me_1'} {\sigma'} \xfrom{\me_2'} \sigma_2$ is given by
    the generated cospan over
    \begin{center}
    \begin{tikzcd}
        {\anyL{\sigma_1}} \ar[dr, dashed] \ar[r, "\anyL{\me_1'}"]
            & {\anyL{\sigma'}}
                & {\anyL{\sigma_2}}
                    \ar[dl, dashed] \ar[l, "\anyL{\me_2'}"'] \\
                & {\img{\anyL{\me_1'}} \setUnion \img{\anyL{\me_2'}}} \ar[u, hook]
    \end{tikzcd}
    \end{center}
    And since it is a local independent unifier we have
    \begin{center}
    \begin{tikzcd}
        {\anyL{\sigma_1}} \ar[dr, dashed] \ar[r, "\ext^{\anyL{\me_1}}\anyL{\me_2}"]
            & {\anyL{\sigma_1} \liplus{\anyL\sigma} \anyL{\sigma_2}}
                & {\anyL{\sigma_2}}
                    \ar[dl, dashed] \ar[l, "\ext^{\anyL{\me_2}}\anyL{\me_1}"'] \\
                & {\img{\anyL{\me_1'}} \setUnion \img{\anyL{\me_2'}}} \ar[u, "\iso"]
    \end{tikzcd}
    \end{center}
    which is exactly the cospan that generates L.I.G.U.
\end{proof}

\begin{lemma}[Pasting Law for Local Independent General Unifiers]\label{lem:ligu-pasting}
    In a diagram
    \begin{tikzcd}
        {\sigma} \ar[r, "\me_1"] \ar[d, "\me_1'"']
            & {\sigma_1} \ar[r, "\me_3"] \ar[d, "\me_2"]
                & {\sigma_3} \ar[d, "\me_3'"]
        \\
        {\sigma_1'} \ar[r, "\me_2'"']
            & {\sigma_2} \ar[r, "\me_3''"']
                & {\sigma_3'}
    \end{tikzcd}
    where the left square is a L.I.G.U., then the whole square is a L.I.G.U. iff the right square is.
\end{lemma}
\begin{proof}
    Immediate from \cref{lem:lic-pasting}.
\end{proof}

\begin{lemma}\label{lem:ligu-rmono-exist}
    Every span of the form $\sigma_1 \xfrom{\me_1} \sigma \xrto{\me_2} \sigma_2$ has a local independent general unifier.
\end{lemma}

\begin{proof}
    We have that the L.I.C. in $\tempC$ commutes
    \begin{tikzcd}[sep=small]
        & {\sigma} 
            \ar[dl, "\tinstU\me_1"'] \ar[dr, -rmonotip, "\tinstU\me_2"] 
    \\ 
    {\sigma_1} 
        \ar[dr, "\ext^{\tinstU\me_1}\tinstU\me_2"']
            && {\sigma_2}
                \ar[dl, "\ext^{\tinstU\me_2}\tinstU\me_1"] 
    \\ 
        & {\sigma_1 \liplus\sigma \sigma_2}
    \end{tikzcd}.
    We need to show that it also commutes in $\tinstC$ with data given by \cref{rem:lic-T-data}.

    Notice that by \cref{lem:lic-vpres}, we have that $\concL{\sigma_1 \liplus\sigma \sigma_2}$ is given by
    \begin{center}
        \begin{tikzcd}
            {\concL\sigma} 
                \ar[r, "\concL{\tinstU\me_2}"] \ar[d, "\concL{\tinstU\me_1}"']
                & {\concL{\sigma_2}}
                    \ar[d, "\concL{\ext^{\tinstU\me_2}\tinstU\me_1}"]
            \\
            {\concL{\sigma_1}}
                \ar[r, "\concL{\ext^{\tinstU\me_1}\tinstU\me_2}"']
                    & {\concL{\sigma_1} \liplus{\concL\sigma} \concL{\sigma_2}}
        \end{tikzcd}
    \end{center}
    and $\ext^{\tinstU\me_1}\tinstU\me_2$ is variable preserving.

    We equip $\ext^{\tinstU\me_1}\tinstU\me_2$ with data $\dheapF(\iso, \anyL{\ext^{\tinstU\me_2}\tinstU\me_1})(\mvarM_{\me_2})$ with the canonical iso
    \[ (\concL{\sigma_1} \liplus{\concL\sigma} \concL{\sigma_2}) \ominus \concL{\ext^{\tinstU\me_1}\tinstU\me_2} \iso \concL{\sigma_2} \ominus \concL{\tinstU\me_2} \]
    We equip $\ext^{\tinstU\me_2}\tinstU\me_1$ with data in the same way.
    This gives that $\mvarM_{\ext^{\tinstU\me_1}\me_2 \comp \me_1} = \mvarM_{\ext^{\tinstU\me_2}\me_1 \comp \me_2}$.
\end{proof}

\subsubsection{Preservation of Connectivity}

Given $\sigma_1 \xsurjto\me \sigma_2 \xto{\me'} \sigma_3$, we ask whether the connectivity of $\me$ is preserved,
that is, whether every location introduced by $\me$ will remain reachable in the composite $\me' \comp \me$.

\begin{definition}[Connectedness Preserving Morphisms]
    A morphism $\me_2 : \sigma_2 \to \sigma_3$ preserves connectivity of a morphism $\me_1 : \sigma_1 \surjto \sigma_2$
    if it satisfies $\anyL{\sigma_2} \pltimes{\anyL{\sigma_3}} (\anyL{\sigma_3} \ominus \anyL{\uncon{\me_2 \comp \me_1}}) = \emptyset$.
\end{definition}

\begin{lemma}[Connectedness Preserving Morphisms Property Equivalence]\label{lem:con-pres-prop-equiv}
    Let $\sigma_1 \xsurjto{\me_1} \sigma_2 \xto{\me_2} \sigma_3$.
    We have the following equivalence.

    \begin{center}
    \begin{tikzcd}
        {\emptyset} 
            \ar[r] \ar[d] \pbsquare{dr} 
                & {\anyL{\sigma_3} \ominus \anyL{\uncon{\me_2 \comp \me_1}}}
                    \ar[d, "\compl{\anyL{\uncon{\me_2 \comp \me_1}}}"] 
        \\
        {\anyL{\sigma_2}} 
            \ar[r, "\anyL{\me_2}"']
                & {\anyL{\sigma_3}}
    \end{tikzcd} $\Leftrightarrow$
    \begin{tikzcd}
        {\anyL{\sigma_2}}
            \ar[d, dashed, "\exists \reach_{\me_1}^{\me_2 \comp \me_1}"'] \ar[rd, "\anyL{\me_2}"]
        \\
        {\anyL{\sigma_{\con{\me_2 \comp \me_1}}}}
            \ar[r, "\anyL{\uncon{\me_2 \comp \me_1}}"' yshift=-5pt] 
            & {\anyL{\sigma_3}}
    \end{tikzcd}
    \end{center}
\end{lemma}
\begin{proof}\nl
    $\Rightarrow$) 
        Consider the canonical iso $\anyL{\sigma_3} \iso \anyL{\sigma_{\con{\me_2 \comp \me_1}}} \oplus (\anyL{\sigma_3} \ominus \anyL{\uncon{\me_2 \comp \me_1}})$.
        Now if any $\va : \vC \in \anyL{\me_2}$ maps to $\anyL{\sigma_{\con{\me_2 \comp \me_1}}}$, we would get a contradiction $\va : \vC \to \emptyset$ and so
        it must factor through $\anyL{\sigma_{\con{\me_2 \comp \me_1}}}$.
    \nl
    $\Leftarrow$)
        We $\img{\anyL{\me_2}} \setInt \img{\compl{\anyL{\uncon{\me_2 \comp \me_1}}}}$ that must be empty as the image of $\anyL{\me_2}$ factors through $\anyL{\sigma_{\con{\me_2 \comp \me_1}}}$. 
\end{proof}

\Cref{lem:con-pres-prop-equiv} tells us that connectivity preservation is exactly uniqueness of factorisation (\cref{lem:fact-uniq-tinstC}),
and so every morphism preserves connectivity as the following lemma shows.

\begin{lemma}[All Morphisms are Connectedness Preserving]
    Every $\me_2 : \sigma_2 \to \sigma_3$ preserves the connectivity of every $\me_1 : \sigma_1 \surjto \sigma_2$.
\end{lemma}
\begin{proof}
    By \cref{lem:fact-uniq-tinstC}, we have
    \begin{center}
        \begin{tikzcd}
            {\sigma_1} 
                \ar[r, twoheadrightarrow, "\me_1"] \ar[drr, twoheadrightarrow, "\con{\me_2 \comp \me_1}"'] \ar[rrr, bend left=30, "\me_2 \comp \me_1"]
                & {\sigma_2}
                    \ar[r, twoheadrightarrow, "\con{\me_2}"]
                    & {\sigma_{\con{\me_2}}}
                        \ar[d, "\iso"] \ar[r, -rmonotip, "\uncon{\me_2}"]
                        & {\sigma_3}
            \\
                    && {\sigma_{\con{\me_2 \comp \me_1}}}
                        \ar[ur, -rmonotip, "\uncon{\me_2 \comp \me_1}"']
        \end{tikzcd}
    \end{center}
    which proves preservation of connectivity by \cref{lem:con-pres-prop-equiv}.
\end{proof}

\subsubsection{Preservation of Disconnectivity}

Similarly, we ask if a morphism $\me' : \sigma_2 \surjto \sigma_3$ preserves a disconnected morphism $\me : \sigma_1 \rto \sigma_2$,
that is, if every location in $\anyL{\sigma_2} \ominus \anyL\me$ remains unreachable in $\me' \comp \me$. 

Such morphisms can only instantiate variables coming from $\sigma_1$ with data that do not point to locations $\anyL{\sigma_2} \ominus \anyL\me$ in order to 
preserve the unreachability of $\anyL{\sigma_2} \ominus \anyL\me$ in the composite $\me_2 \comp \me_1$.

\begin{definition}[Disconnectedness Preserving Morphisms]\label{def:dcon-pres-morph}
    A morphism $\me_2 : \sigma_2 \to \sigma_3$ preserves the disconnectivity of a morphism $\me_1 : \sigma_1 \rto \sigma_2$
    if it satisfies $(\anyL{\sigma_2} \ominus \anyL{\me_1}) \pltimes{\anyL{\sigma_3}} \anyL{\sigma_{\con{\me_2 \comp \me_1}}} = \emptyset$.
\end{definition}

\begin{lemma}[Disconnectedness Preserving Morphisms Property Equivalence]\label{lem:dcon-pres-prop-equiv}
    Let $\sigma_1 \xrto{\me_1} \sigma_2 \xto{\me_2} \sigma_3$.
    We have the following equivalence.

    \begin{tikzcd}[sep=small]
        {\emptyset} 
            \ar[r] \ar[dd] \pbsquare{ddr} 
            & {\anyL{\sigma_2} \ominus \anyL{\me_1}}
                \ar[d, "\compl{\anyL{\me_1}}"] 
        \\
            & {\anyL{\sigma_2}}
                \ar[d, "\anyL{\me_2}"]
        \\
        {\anyL{\sigma_{\con{\me_2 \comp \me_1}}}} 
            \ar[r, -rmonotip, "\anyL{\uncon{\me_2 \comp \me_1}}"' yshift=-3pt] 
            & {\anyL{\sigma_3}}
    \end{tikzcd}$\mathclap{\xLeftrightarrow{(1)}}$
    \begin{tikzcd}
        {\anyL{\sigma_1}}
            \ar[d, "\anyL{\con{\me_2 \comp \me_1}}"] \ar[r, -rmonotip, "\anyL{\me_1}"] \pbsquare{dr}
            & {\anyL{\sigma_2}}
                \ar[d, "\anyL{\me_2}"'] 
        \\
        {\anyL{\sigma_{\con{\me_2 \comp \me_1}}}}
            \ar[r, -rmonotip, "\anyL{\uncon{\me_2 \comp \me_1}}"' yshift=-5pt] 
            & {\anyL{\sigma_3}}
    \end{tikzcd}$\hspace{5pt}\mathclap{\xLeftrightarrow{(2)}}$
    \begin{tikzcd}
        {\sigma_1}
            \ar[d, twoheadrightarrow, "\con{\me_2 \comp \me_1}"] \ar[r, -rmonotip, "\me_1"] \pbsquare{dr}
            & {\sigma_2}
                \ar[d, "\me_2"'] 
        \\
        {\sigma_{\con{\me_2 \comp \me_1}}}
            \ar[r, -rmonotip, "\uncon{\me_2 \comp \me_1}"' yshift=-5pt] 
            & {\sigma_3}
    \end{tikzcd}$\hspace{5pt}\mathclap{\xLeftrightarrow{(3)}}$
    \begin{tikzcd}[sep=small]
        {\anyL{\sigma_2} \ominus \anyL{\me_1}}
            \ar[d, dashed, "\exists \anyL{\me_2}\res|{\me_1}"'] \ar[r, "\compl{\anyL{\me_1}}"] 
            & {\anyL{\sigma_2}}
                \ar[d, "\anyL{\me_2}"'] 
        \\
        {\anyL{\sigma_3} \ominus \anyL{\uncon{\me_2 \comp \me_1}}}
            \ar[r, "\compl{\anyL{\uncon{\me_2 \comp \me_1}}}"' yshift=-5pt] 
            & {\anyL{\sigma_3}}
    \end{tikzcd}
\end{lemma}
\begin{proof}
    For (1), we consider the canonical isomorphism $\anyL{\sigma_2} \iso \anyL{\sigma_1} \oplus (\anyL{\sigma_2} \ominus \anyL{\me_1})$.
    For $\xRightarrow{(1)}$ we cannot have anything factoring through $(\anyL{\sigma_2} \ominus \anyL{\me_1})$ as we get a contradiction 
    and for $\xLeftarrow{(1)}$, if $\anyL{\sigma_{\con{\me_2 \comp \me_1}}} \xfrom{\pi_1} \vw \xto{\pi_2} (\anyL{\sigma_2} \ominus \anyL{\me_1})$ is the pullback
    on the left, we can form a square in the right one to get a unique injection $\anyL{\sigma_{\con{\me_1}}} \oplus \vw \dashto \anyL{\sigma_{\con{\me_1}}}$
    and so $\vw$ must be empty.
    For $\xRightarrow{(2)}$, by \cref{lem:pullback-vpres}, the diagram in $\tempC$ is a pullback and by 
    \cref{lem:tinstU-pls-refl} it is also a pullback in $\tinstC$.
    For $\xLeftarrow{(3)}$, by \cref{lem:anyL-pls-pres} the $\anyL\dash$-image is a pullback.

    For $\xLeftarrow{(3)}$, we have that 
    \begin{tikzcd}[sep=small]
        {\emptyset} 
            \ar[r] \ar[dd]
            & {\anyL{\sigma_2} \ominus \anyL{\me_1}}
                \ar[r, dashed, "\anyL{\me_2}\res|{\me_1}" yshift=3pt] \ar[d, "\compl{\anyL{\me_1}}"']
                & {\anyL{\sigma_3} \ominus \anyL{\uncon{\me_2 \comp \me_1}}}
                    \ar[ddl, "\compl{\anyL{\uncon{\me_2 \comp \me_1}}}"]
        \\
            & {\anyL{\sigma_2}}
                \ar[d, "\anyL{\me_2}"']
        \\
        {\anyL{\sigma_{\con{\me_2 \comp \me_1}}}}
            \ar[r, "\anyL{\uncon{\me_2 \comp \me_1}}"' yshift=-3pt]
            & {\anyL{\sigma_3}}
    \end{tikzcd}
    is a pullback as $\img{\anyL{\uncon{\me_2 \comp \me_1}}} \setInt \img{\compl{\anyL{\uncon{\me_2 \comp \me_1}}}} = \emptyset$ and so the inner square must be a pullback.

    For $\xRightarrow{(3)}$, the second pullback gives that $\anyL{\me_2}(\compl{\anyL{\me_1}}(\va)) \not\in \img{\anyL{\uncon{\me_2 \comp \me_1}}}$.
    
\end{proof}

\begin{lemma}[Regular Monos are Disconnectedness Preserving]\label{lem:dcon-pres-rmono}
    For $\sigma_1 \xrto{\me_1} \sigma_2 \xrto{\me_2} \sigma_3$, $\me_2$ preserves disconnectivity of $\me_1$.
\end{lemma}
\begin{proof}
    Notice that $\con{\me_2 \comp \me_1}$ is epi and a regular mono and thus iso.
    As a result,
    \begin{center}
        \begin{tikzcd}
            {\anyL{\sigma_1}}
            \ar[r, -rmonotip, "\anyL{\me_1}"] \ar[d, "\anyL{\con{\me_2 \comp \me_1}}"', "\iso"] \pbsquare{dr}
                    & {\anyL{\sigma_2}}
                        \ar[d, -rmonotip, "\anyL{\me_2}"]
            \\
            {\anyL{\sigma_{\con{\me_2 \comp \me_1}}}}
                \ar[r, -rmonotip, "\anyL{\uncon{\me_2 \comp \me_1}}"' yshift=-3pt]
                    & {\anyL{\sigma_3}}
        \end{tikzcd}
    \end{center}
    is a pullback by \cref{lem:pullback-iso-mono}.
\end{proof}

\begin{lemma}[Disconnectedness Preservation is Downwards Closed]\label{lem:dcon-pres-dclosed}
    For $\sigma_1 \xrto{\me_1} \sigma_2 \xto{\me_2} \sigma_3 \xto{\me_3} \sigma_4$, if $\me_3 \comp \me_2$ preserves disconnectivity of $\me_1$ then
    $\me_2$ preserves disconnectivity of $\me_1$.
\end{lemma}
\begin{proof}
    We have a diagram
    \begin{center}
        \begin{tikzcd}
        {\emptyset}
            \ar[r] \ar[dd]
            & {\anyL{\sigma_2} \ominus \anyL{\me_1}} 
                \ar[d, "\compl{\anyL{\me_1}}"]
        \\
            & {\anyL{\sigma_2}}
                \ar[d, "\anyL{\me_2}"]
        \\
        {\anyL{\sigma_{\con{\me_2 \comp \me_1}}}}
            \ar[r, "\anyL{\uncon{\me_2 \comp \me_1}}"] \ar[d, dashed, "\anyL{\lt{\me_3 \comp \uncon{\me_2 \comp \me_1}}}"']
            & {\anyL{\sigma_3}}
                \ar[d, "\anyL{\me_3}"]
        \\
        {\anyL{\sigma_{\con{\me_3 \comp \me_2 \comp \me_1}}}}
            \ar[r, "\anyL{\uncon{\me_3 \comp \me_2 \comp \me_1}}"' yshift=-10pt]
            & {\anyL{\sigma_4}}
        \end{tikzcd}
    \end{center}
    Where the dashed morphism is given by \cref{lem:ofs-tinstC}.
    That $\me_3 \comp \me_2$ preserves disconnectivity of $\me_1$ gives us that the outer square is a pullback,
    and by \cref{lem:pullback-initial-obj}, the inner square is also pullback which translates to 
    $\me_2$ preserves disconnectivity of $\me_1$.
\end{proof}

\begin{lemma}[L.I.G.U. and Disconnectedness Preservation]\label{lem:ligu-dcon}
    Every L.I.G.U. diagram of the form 
    \begin{tikzcd}
        & {\sigma_1} 
            \ar[dl, "\me_1'"'] \ar[dr, -rmonotip, "\me_1"]
    \\
    {\sigma_2'}
        \ar[dr, "\me_2'"']
            && {\sigma_2}
                    \ar[dl, "\me_2"]
    \\
        & {\sigma_3}
    \end{tikzcd}
    satisfies $\me_2$ preserves the disconnectivity of $\me_1$.
\end{lemma}
\begin{proof}
    By \cref{lem:pullback-vpres} we have that the $\tinstU$-image is a pullback in $\tempC$,
    which by \cref{lem:tinstU-pls-refl} it is a pullback in $\tinstC$.

    So we have a pullback
    \begin{center}
        \begin{tikzcd}
            {\sigma_1}
                \ar[r, "\me_1'"] \ar[d, -rmonotip, "\me_1"'] \pbsquare{dr}
                    & {\sigma_2'}
                        \ar[d, -rmonotip, "\me_2'"]
            \\
            {\sigma_2} \ar[r, "\me_2"']
                    & {\sigma_3}
        \end{tikzcd}
    \end{center}
    and notice $\me_2'$ is r.mono by \cref{lem:lic-vpres} and $\tinstU$ reflecting variable preserving morphisms. 
    
    By \cref{lem:pullback-factoring}, we have that 
    \begin{center}
        \begin{tikzcd}
            {\sigma_1}
                \ar[r, twoheadrightarrow, "\con{\me_1'}"] \ar[dd, -rmonotip, "\me_1"'] \pbsquare{ddr}
                    & {\sigma_{\con{\me_2' \comp \me_1'}}}
                        \ar[d, -rmonotip, "\uncon{\me_1'}"]
            \\
                    & {\sigma_2} 
                        \ar[d, -rmonotip, "\me_2'"]
            \\
            {\sigma_2} \ar[r, "\me_2"']
                    & {\sigma_3}
        \end{tikzcd}
    \end{center}
    which proves that $\me_2$ preserves disconnectivity of $\me_1$ by \cref{lem:dcon-pres-prop-equiv}. 
\end{proof}

\begin{lemma}[Disconnectedness Preserving Morphisms are Data Agnostic]\label{lem:dcon-pres-data-agn}
    For 
    \begin{tikzcd}[sep=small]
        {\sigma_1} 
            \ar[r, -rmonotip, "\me_1"]
            & {\sigma_2} \ar[r, "\me_2"]
                & {\sigma_3}
    \end{tikzcd}
    if $\me_2$ preserves disconnectivity of $\me_1$, then for all $\mvarM' : \dheapF(\concL{\sigma_2} \ominus \concL{\me_1}, \anyL{\sigma_2})$,  $\me_2$ preserves disconnectivity of $\me_1' = \pair{\instU\me_1}{\mvarM'}$.
\end{lemma}
\begin{proof}
    We first show that $\anyL{\sigma_1} \xto{\anyL{\con{\me_2 \comp \me_1}}} \anyL{\sigma_{\con{\me_2 \comp \me_1}}} \xto{\anyL{\uncon{\me_2 \comp \me_1}}} \anyL{\sigma_3}$
    satisfies \cref{def:reach-loc} for $\mvarM_{\me_2 \comp \me_1'}$.
    This would give us a $\vh$ that satisfies
    \begin{center}
        \begin{tikzcd}
            {\anyL{\sigma_1}} 
                \ar[r, "\anyL{\con{\me_2 \comp \me_1}}"] \ar[dr, "\anyL{\con{\me_2 \comp \me_1'}}"']
                & {\anyL{\sigma_{\con{\me_2 \comp \me_1}}}}
                    \ar[r, "\uncon{\me_2 \comp \me_1}"]
                        & {\anyL{\sigma_3}}
            \\
                & {\anyL{\sigma_{\con{\me_2 \comp \me_1'}}}}
                    \ar[ur, "\uncon{\me_2 \comp \me_1'}"'] \ar[u, dashed, "\vh"]
        \end{tikzcd}
    \end{center}
    Property (i) follows immediately.
    For Property (ii), first by the lifting property we have
    \begin{tikzcd}
        {\sigma_{\con{\me_2 \comp \me_1}}} 
            \ar[rr, -rmonotip, "\uncon{\me_2 \comp \me_1}"] \ar[dr, dashed, -rmonotip, "\lt{\uncon{\me_2 \comp \me_1}}"']
                && {\sigma_3}
        \\
            & {\sigma_{\con{\me_2}}}
                \ar[ur, -rmonotip, "\uncon{\me_2}"']
    \end{tikzcd},
    and the pullback $\anyL{\sigma_{\con{\me_2 \comp \me_1}}} \pltimes{\anyL{\sigma_3}} (\concL{\sigma_3} \ominus \concL{\me_2 \comp \me_1})$
    is given by the pullback
    ${\concL{\sigma_{\con{\me_2 \comp \me_1}}} \pltimes{\concL{\sigma_{\con{\me_2}}}} (\concL{\sigma_{\con{\me_2}}} \ominus \concL{\con{\me_2} \comp \me_1})}$
    from
    \begin{center}
        \begin{tikzcd}
            {\anyL{\sigma_{\con{\me_2 \comp \me_1}}}}
                \ar[r, "\anyL{\lt{\uncon{\me_2 \comp \me_1}}}"] \ar[rr, bend left=15, "\anyL{\uncon{\me_2 \comp \me_1}}"] 
                    & {\anyL{\sigma_{\me_2}}}
                        \ar[r, "\anyL{\uncon{\me_2}}"]
                        & {\anyL{\sigma_3}}
            \\
            {\concL{\sigma_{\con{\me_2 \comp \me_1}}}}
                \ar[r, "\concL{\lt{\uncon{\me_2 \comp \me_1}}}"] \ar[u, "\concL{\sigma_{\con{\me_2 \comp \me_1}}}"] \pbsquare{ur}
                    & {\concL{\sigma_{\con{\me_2}}}}
                        \ar[r, "\concL{\uncon{\me_2}}"] \ar[u, "\concL{\sigma_{\con{\me_2}}}"] \pbsquare{ur}
                        & {\concL{\sigma_3}}
                            \ar[u, "\concL{\sigma_3}"']
            \\
            {\sigma_{\concL{\con{\me_2 \comp \me_1}}}
             \pltimes\relax
             (\concL{\sigma_{\con{\me_2}}} \ominus \concL{\con{\me_2} \comp \me_1})}
                \ar[r] \ar[u] \pbsquare{ur}
                & {\concL{\sigma_{\con{\me_2}}} \ominus \concL{\con{\me_2} \comp \me_1}}
                    \ar[r, "\concL{\uncon{\me_2}}"' yshift=-4pt] \ar[u, "\compl{\concL{\con{\me_2} \comp \me_1}}"] \pbsquare{ur}
                    & {\concL{\sigma_3} \ominus \concL{\me_2 \comp \me_1}}
                        \ar[u, "\compl{\concL{\me_2 \comp \me_1}}"']
        \end{tikzcd}
    \end{center}
    Now by \cref{lem:dcon-pres-dclosed}, we have that $\con{\me_2}$ preserves disconnectivity of $\me_1$ and so
    \begin{center}
    \begin{tikzcd}
        {\anyL{\sigma_1}}
            \ar[r, "\anyL{\me_1}"] \ar[d, "\anyL{\con{\me_2 \comp \me_1}}"'] \pbsquare{dr}
                & {\anyL{\sigma_2}}
                    \ar[d, "\anyL{\con{\me_2}}"]
        \\
        {\anyL{\sigma_{\con{\me_2 \comp \me_1}}}}
            \ar[r, "\anyL{\lt{\uncon{\me_2 \comp \me_1}}}"']
                & {\anyL{\sigma_{\con{\me_2}}}}
    \end{tikzcd} 
    \end{center}
    which implies that $\res_{\compl{\concL{\con{\me_2} \comp \me_1}}}\concL{\lt{\uncon{\me_2 \comp \me_1}}}$ must factor as
    \begin{center}
    \begin{tikzcd}
            {\concL{\sigma_{\con{\me_2 \comp \me_1}}}
             \pltimes\relax
             (\concL{\sigma_{\con{\me_2}}} \ominus \concL{\con{\me_2} \comp \me_1})}
                \ar[rr, "\res_{\compl{\concL{\con{\me_2} \comp \me_1}}}\concL{\lt{\uncon{\me_2 \comp \me_1}}}"] \ar[dr, dashed]
                    && {\concL{\sigma_{\con{\me_2}}} \ominus \concL{\con{\me_2} \comp \me_1}}
            \\
                & {\concL{\sigma_{\con{\me_2}}} \ominus \concL{\con{\me_2}}}
                    \ar[ur, hook]
    \end{tikzcd}
    \end{center}
    As a result, all locations $\anyL{\sigma_{\con{\me_2 \comp \me_1}}} \pltimes{\anyL{\sigma_3}} (\concL{\sigma_3} \ominus \concL{\me_2 \comp \me_1})$
    in $\mvarM_{\me_2 \comp \me_1'}$ point to cells in $\mvarM_{\con{\me_2}}$.
    Since $\tinstU\me_1 = \tinstU\me_1'$, the locations in these cells factor through $\anyL{\sigma_{\con{\me_2 \comp \me_1}}}$ by \cref{def:reach-loc}
    for $\anyL{\con{\me_2 \comp \me_1}}$ at $\mvarM_{\me_2 \comp \me_1}$.

    We now use $\vh$ to show
    \begin{center}
        \begin{tikzcd}
            {\anyL{\sigma_1}} 
                \ar[r, Rightarrow] \ar[d, "\anyL{\con{\me_2 \comp \me_1'}}"'] \pbsquare{dr}
                & {\anyL{\sigma_1}}
                    \ar[d, "\anyL{\con{\me_2 \comp \me_1}}"'] \ar[rr, "\anyL{\me_1}"] \pbsquare{drr}
                        && {\anyL{\sigma_2}} 
                            \ar[d, "\anyL{\me_2}"]
            \\
            {\anyL{\sigma_{\con{\me_2 \comp \me_1'}}}}
                \ar[r, dashed, "\vh"] \ar[rrr, bend right=30, "\anyL{\uncon{\me_2 \comp \me_1'}}"']
                & {\anyL{\sigma_{\con{\me_2 \comp \me_1}}}}
                    \ar[rr, "\anyL{\uncon{\me_2 \comp \me_1}}"]
                        && {\anyL{\sigma_3}}
        \end{tikzcd}
    \end{center}
    The right square is a pullback by \cref{lem:dcon-pres-prop-equiv}, and the left square is a pullback by \cref{lem:pullback-iso-mono}. 
    As a result by pullback pasting the outer square is a pullback giving us that $\me_2$ preserves the disconnectivity of $\me_1'$.
\end{proof}

\Cref{lem:dcon-pres-data-agn} leads us to a definition of disconnectivity preservation w.r.t. morphisms in $\tempC$.

\begin{definition}[Disconnectedness Preserving Morphisms w.r.t. $\tempC$]
    We say that $\me_2 : \sigma_2 \to \sigma_3$ preserves the disconnectivity of $\sigma_1 \xrto{\rho_1} \sigma_2$ living in $\tempC$, 
    if for all $\mvarM_{\rho_1} : \dheapF(\concL{\sigma_2} \ominus \concL{\rho_1}, \anyL{\sigma_2})$,
        $\me_2$ preserves the disconnectivity of $\pair{\rho_1}{\mvarM_{\rho_1}}$.
\end{definition}

\begin{lemma}[Disconnectedness Preservation and L.I.C.]
    For a span $\sigma_1 \xfrom{\tinstU\me_1} \sigma \xrto{\rho_2} \sigma_2$ in $\tempC$,
    the induced morphism in $\tinstC$, $\ext^{\rho_2}{\me_1}$, given by \cref{lem:lic-exist} and \cref{rem:lic-T-data},
    satisfies: $\ext^{\rho_2}{\me_1}$ preserves the disconnectivity of $\rho_2$.
\end{lemma}
\begin{proof}
    For some $\mvarM_{\rho_2} : \dheapF(\concL{\sigma_2} \ominus \concL{\rho_2}, \anyL{\sigma_2})$, and $\me_2 = \pair{\rho_2}{\mvarM_{\rho_2}}$,
    we have by \cref{lem:ligu-rmono-exist} that
    \begin{tikzcd}
        & {\sigma}
            \ar[dl, "\me_1"'] \ar[dr, -rmonotip, "\me_2"]
    \\
    {\sigma_1}
        \ar[dr, "\ext^{\tinstU\me_1}\me_2"']
            && {\sigma_2}
                \ar[dl, "\ext^{\tinstU\me_2}\me_1"]
    \\
        & {\sigma_1 \liplus\sigma \sigma_2}
    \end{tikzcd}
    with data given by \cref{rem:lic-T-data} commutes in $\tinstC$. 
    Since this is a L.I.G.U. we have by \cref{lem:ligu-dcon} that $\ext^{\tinstU\me_2}\me_1$ preserves the disconnectivity of $\me_2$.
\end{proof}
